\section{Equivariant \texorpdfstring{$K$}{K}-theory ring of the affine Grassmannian}\label{section: equivariant K-theory ring of the affine grassmannian}

The aim of this section is to establish an algebro-geometric description of the ring $K^{\top, 0}_{T \times \Gm^{\rot}}(\Gr_{G}; \C)$, the main focus of this paper. Our Theorem \ref{theorem: equivariant K-theory as deformation to normal cone -- G sssc} presents this ring in terms of completed global functions on affine schemes of representation-theoretical origin. These are given by gluing together infinitely many deformations to normal cones at quasi-diagonals $Z_{\ell}$ in two copies of the torus $T$. We will see that this description can be made very explicit.


Here, let us point out some related literature. The historically first description \cite{KK90} gives $K$-theory of the affine Grassmannian as a dual to its compactly supported variant. To further describe the ring structure, one can appeal to the GKM description \cite{HHH05}. However, this only realizes $K^{\top, 0}_{T \times \Gm^{\rot}}(\Gr_G; \C)$ as a subring of an infinite product of Laurent polynomial rings, cut out by an infinite set of equalizer conditions -- to get a concrete understanding of the ring, one needs to solve this nontrivial system of congruences. Our result may be interpreted as indirectly doing so.


For most of the discussion, we assume $G=G^{\sc}$ to be simply connected or $GL_n$. (While the definitions work for any reductive $G$ as stated, they have to be carefully modified to give the correct answer for $K$-theory.)

We first recall the case of cohomology in \S \ref{section: recollaction of cohomology} and then the GKM description of $K$-theory in \S \ref{section: Goresky--Kottwitz--McPherson description of equivariant $K$-theory}.
We then move on to define the quasi-diagonals inside $T \times \Gm^{\rot} \times T \sslash W$ in \S \ref{section: quasi-diagonals at roots of unity} and compare them to specializations of cocharacter graphs. We blow up these quasi-diagonals in \S \ref{section: deformations of quasi-diagonals to normal cones}, obtaining the deformation to their normal cones. We then explain in \S \ref{section: map to equivariant K-theory: combinatorics}, in purely combinatorial terms, how the completed ring of functions on this deformation maps to the GKM ring of the affine Grassmannian. We then recast this argument in terms of classes of vector bundles in \S \ref{section: map to equivariant K-theory: topology}. In fact, one may view these arguments independently. Moreover, both constructions work integrally. We then prove our main result, Theorem \ref{theorem: equivariant K-theory as deformation to normal cone -- G sssc}, with $\C$-coefficients for simply connected group $G$ in \S \ref{section: bezrukavnikov--finkelberg description of equivariant K-theory}.
%
We discuss its generalizations in \S \ref{section: variants and generalizations} -- this includes arbitrary reductive groups, arbitrary affine flag varieties, and change of equivariance.


\subsection{Recollection of cohomology}\label{section: recollaction of cohomology}
Let $k=\C$. Before we embark to describe the $K$-theory, we recall the well-known case of cohomology of $\Gr_G$, following \cite{Gin95, BF07, Yun10}. This argument has two parts: the purely topological computation of non-equivariant cohomology of $\Gr_G$, and the subsequent description of its equivariant version as a deformation to normal cone.

\subsubsection{Cohomology of finite flag varieties.}
As a warmup, recall the equivariant cohomology of finite flag varieties, which will be useful later. Let $\k$ be a coefficient ring.
\begin{lemma}\label{lemma: equivariant cohomology of finite flag varieties}
Let $G$ be a reductive group and $T \leq G$ be a maximal torus over $k$. Let $P_J \leq G$ be a parabolic subgroup and $W_2 \leq W$ corresponding Weyl group. Assume $\k$ is a field of $\chara \k = 0$. Then we have a graded ring isomorphism
\begin{equation*}
    \O(\t \times_{\t \sslash W} \t \sslash W_2) \xrightarrow{=} H^{\bullet}_{T}(G/P_J; \k),
\end{equation*}
where the grading on the left-hand side comes from the weight two action on both copies of $\t$.
More generally, take a pair of Levi subgroups $L_1, L_2 \leq G$ with Weyl groups $W_1, W_2 \leq W$. Denote $P_J$ the standard parabolic associated to $L_J := L_2$. Then we have a graded ring isomorphism
\begin{equation*}
    \O(\t \sslash W_1 \times_{\t \sslash W} \t \sslash W_2) \xrightarrow{=} H^{\bullet}_{L_1}(G/P_J; \k).
\end{equation*}
\end{lemma}
\begin{proof}
The $G$-equivariant cohomology is given by 
\begin{equation*}
 H^{\bullet}_G(G/P_J; \k) = H^{\bullet}(G \backslash G / P_J; \k) = H^{\bullet}(\pt / P_J; \k) = H^{\bullet}(\pt / L_J; \k) = H^{\bullet}_{L_J}(\pt; \k),   
\end{equation*}
the equivariant base for $L_J$. Consequently, we have a natural identification
\begin{equation*}
 \O(\t \sslash W_J) \xrightarrow{=} H^{\bullet}_G(G/P_J; \k)   
\end{equation*}
compatibly with the structure map to the equivariant base $\O(\t \sslash W) \xrightarrow{=} H^{\bullet}_G(\pt; \k)$. By equivariant formality, the $T$-equivariant cohomology is given by base change along $H^{\bullet}_G(\pt; \k) \to H^{\bullet}_T(\pt; \k)$, so we get the desired
\begin{equation*}
\O(\t \times_{\t \sslash W} \t \sslash W_J) \xrightarrow{=} H^{\bullet}_{T}(G/P_J; \k).    
\end{equation*}
Similarly, $L_1$-equivariant cohomology is given by base change along $H^{\bullet}_G(\pt; \k) \to H^{\bullet}_{L_1}(\pt; \k)$, proving the more general statement.
\end{proof}



\begin{remark}\label{remark: finite flag varieties with trivial loop rotation}
Consequently, adding the equivariance with respect to the trivial loop rotation, we have
\begin{equation*}
   \O(\Ga^{\rot} \times \t \sslash W_1 \times_{\t \sslash W} \t \sslash W_2) \xrightarrow{=} H^{\bullet}_{L_1 \times \Gm^{\rot}}(G/P_J; \k). 
\end{equation*}
In particular, this statement holds after completion $(-)^{\wedge}_{(0, 0)}$ over the origin $(0,0) \in \t \times \Ga^{\rot}$.
\end{remark}

\begin{remark}\label{remark: GKM for cohomology of finite flag varieties}
In terms of their GKM presentations, this description works as follows. Take the diagonal $\t\sslash W \hookrightarrow \t \sslash W \times \t \sslash W$. Consider its base change $\Delta^J := \t \times_{\t\sslash W} \t\sslash W_J \hookrightarrow \t \times \t\sslash W_J$. This defines a closed subscheme parametrizing $\{ (x, y) \in \t \times \t \sslash W_J \mid x = y \mod W  \}$. It has irreducible components isomorphic to copies of $\t$ labeled by the finite set $W/W_J$. The GKM localization to fixed-points corresponds to restriction to these components. 
\end{remark}


\begin{remark}\label{remark: cohomology of finite flag varietes as coinvariant algebra}
Taking fiber over the origin in the equivariant base, we get the usual description of non-equivariant cohomology, namely
\begin{equation*}
    H^{\bullet}(G/P_J; \k) = \frac{\k[t_1, \dots, t_n]}{\k[t_1, \dots, t_n]^W}. 
\end{equation*}
\end{remark}


\begin{remark}
Upon taking the correct equivariant base, the discussion in this section works for any equivariant multiplicative generalized cohomology theory $E_G$. In particular, it works for $K$-theory: for any coefficient ring $\k$, for example $\k=\Z$, we have
\begin{equation*}
    \O(T \times_{T \sslash W} T \sslash W_J) = K^{\top, 0}_{T \times \Gm^{\rot}}(G/P_J; \k).
\end{equation*}
All the above remarks are valid in this setting.
\end{remark}









\subsubsection{Based loop space description of ordinary cohomology.}\label{section: based loop space description}
Assume $G$ is a simply connected group over $k =\C$. The non-equivariant cohomology of the affine Grassmannian $\Gr_G$ can be computed by the following classical argument, going back to \cite[\S 1.7]{Gin95}. 

Firstly, denote $K$ a connected compact real Lie group with complexification $G(\C)$. Then $K \leq G(\C)$ is a maximal compact subgroup. Denote $\Omega K$ its based loop space. It is then a nontrivial fact that we can identify the homotopy type of the affine Grassmannian as $\Gr_G(\C) \simeq \Omega K$, see for example \cite[Introduction]{YZ09}, \cite[\S 1.2]{Gin95}, \cite[Remark 5.1]{HHH05}, \cite[\S 1.6]{Zhu15}.

We are thus left to understand the graded ring $H^{\bullet}(\Omega K)$. Firstly, the cohomology $H^{\bullet}(K)$ is understood by \cite{Bor53}. In particular, with coefficients in a field $\k$ of $\chara \k = 0$, the cohomology ring $H^{\bullet}(K; \k) = \Lambda_{\k}[x_1, \dots, x_n]$ is the exterior algebra on finitely many generators $x_i$, $i=1, \dots, d$ by \cite[Proposition 7.2]{Bor53}. Furthermore, $d = \rank G$ and $\deg x_i = a_i -1$ where $a_1, \dots a_d$ are the exponents of $G$.

Consequently, the rational homotopy type of $K$ is that of a product of odd dimensional spheres $S^{a_i-1}$. By the compatibility of the based loop space $\Omega(-)$ with Eilenberg-MacLane spaces, the rational homotopy type of $\Omega K$ is given by a product of even dimensional spheres $S^{a_i - 2}$, and we deduce
\begin{equation}\label{equation: cohomology of loop groups}
    H^{\bullet}(\Omega K; \k) = \k[y_1, \dots, y_n], \qquad \deg y_i = a_i - 2
\end{equation}
In a more hands-on way, the passage from $K$ to $\Omega K$ can be also done directly through the path space fibration $\Omega K \to P K \to K$ and Serre's spectral sequence; the elements $x_i$ are transgressions of $y_i$. See \cite[Example 1.16]{Hat04}. 




\subsubsection{Bezrukavnikov--Finkelberg description of equivariant cohomology.}
We recall the description of the equivariant cohomology of affine flag varieties over $k=\C$ via certain deformations to normal cones, following \cite{BF07, Yun10}. This crucially uses the topological input \S \ref{section: based loop space description}. We assume the coefficient ring $\k$ is a field of $\chara \k = 0$.

Let us sketch this argument.
The equivariant cohomology of $\Fl^{J} = \lop G / \eP_{J}$ is a module over the equivariant base via
\begin{equation}\label{equation: equivariant cohomology as modules over equivariant base}
    H^{\bullet}_{G \times \Gm^{\rot} \times L_{J}}(\pt; \k) \to H^{\bullet}_{\lop G \times \Gm^{\rot}}(\Fl^{J}; \k).
\end{equation}
The map \eqref{equation: equivariant cohomology as modules over equivariant base} is easily seen to be injective. The source is given by space $\O(\t \sslash W \times \Ga^{\rot} \times \t \sslash W_{L_J})$. In fact, both sides are abstractly isomorphic polynomial rings.

However, \eqref{equation: equivariant cohomology as modules over equivariant base} is not surjective. To correct for the missing functions, one needs to deform to the normal cone of the image of $\t$  under the diagonal projection $\t \to \t \sslash W \times \t \sslash W_{L_J}$, in other words to the normal cone of the closed subscheme $\t \sslash W \times_{\t \sslash W} \t \sslash W_{L_J}$.
One directly checks that \eqref{equation: equivariant cohomology as modules over equivariant base} factors injectively through this deformation.
Finally, to prove surjectivity, one notes that \eqref{equation: equivariant cohomology as modules over equivariant base} is a graded map with respect to the natural gradings on both sides. The graded pieces are finite-dimensional, so it is enough to compare their Hilbert series. These are easily computable from the deformation description for the left-hand side, and \eqref{equation: cohomology of loop groups} together with equivariant formality for the right-hand side.


To state the result, consider Levi subgroups $L_1, L_2$ of $G$, with associated Weyl groups $W_1, W_2 \leq W$. 
Let $\eP_J$ be the standard parahoric sugroup associated to $L_J :=L_2$. Conisder the corresponding affine flag variety $\Fl^J = \lo G / \eP_J$, with the natural left action $\lop L_1 \rtimes \Gm^{\rot} \acts \Fl^J$.
\begin{proposition}\label{proposition: cohomology of affine flag varieties as deformation to the normal cone}
 The $(\lop L_1 \times \Gm^{\rot})$-equivariant cohomology ring of the affine flag variety $\Fl^{J}$ is given, together with its structure map, by  
 \begin{equation}
    \begin{tikzcd}
        \O(N( \t \sslash W_1 \times_{\t \sslash W} \t \sslash W_2, \t \sslash W_1 \times \t \sslash W_2)) \arrow[r, equal]& H_{\lop L_1 \times \Gm^{\rot}}(\Fl^{J}; \k) \\
         \O(\t \sslash W_1 \times \Ga^{\rot} \times \t \sslash W_2) \arrow[u, hookrightarrow] \arrow[r, equal] & H_{L_1 \times \Gm^{\rot} \times L_{2}}(\pt; \k) \arrow[u, hookrightarrow] 
    \end{tikzcd}
 \end{equation}
\end{proposition}
\begin{proof}
The case of $\lop G \rtimes \Gm^{\rot} \acts \Gr_G$ is covered in \cite[Theorem 1]{BF07}. From there, the general case follows by the argument \cite[Theorem 4.3]{Yun10}. Strictly speaking, \cite{Yun10} treats only the full flag variety with the Iwahori action, but the argument works in general.
\end{proof}

\begin{remark}[Informal GKM localization]\label{remark: informal GKM localization}
Working equivariantly with respect to $T \times \Gm^{\rot}$, the above identification can be unraveled in GKM terms. Let us give an informal account in the complex analytic setting now; we will give a detailed discussion in the $K$-theoretical context below.

On the complex analytic level, consider $\t(\C) \sslash W_{\aff} \hookrightarrow \t(\C) \sslash W_{\aff} \times \Ga \times \t(\C) \sslash W_{\aff}$. By pullback, we obtain the subspace $\overline{\mathbf{\Gamma}}(\C) \hookrightarrow \t(\C) \times \Ga(\C) \times \t(\C) \sslash W_{\aff}$. Its irreducible components are copies of $\t(\C)$, labeled by the cocharacter lattice $X_{\bullet}$. Take the strict transform $\overline{\mathbf{\Upsilon}}(\C)$ of $\overline{\mathbf{\Gamma}}(\C)$ in the blowup $N( \t(\C) \sslash W_1 \times_{\t(\C) \sslash W} \t(\C) \sslash W_2, \t(\C) \sslash W_1 \times \t(\C) \sslash W_2)$. Its irreducible components are still copies of $\t(\C)$ labeled by $X_{\bullet}$, and the localization corresponds to restriction of functions to them. Also see \cite[\S 3.2]{BF07}.
\end{remark}




\subsection{Goresky--Kottwitz--McPherson description of equivariant \texorpdfstring{$K$}{K}-theory} \label{section: Goresky--Kottwitz--McPherson description of equivariant $K$-theory}
We record some standard formality results on equivariant topological $K$-theory of $\Gr_G$, together with its well-known GKM description. In this subsection, $G$ is any reductive group over $k=\C$.


\subsubsection{Limit along affine Schubert varieties.} Recalling \S \ref{subsubsection: K-theory of colimits}, consider the affine Grassmannian $\Gr_G$ with its $T\times \Gm^{\rot}$-equivariant Iwahori paving by affine cells. Consider the induced skeleta $\sk_i = \sk_i \Gr$, $i \in \N_0$. By induction, $K^{\top, \bullet}_{T \times \Gm^{\rot}}(\sk_i)$ is a free module over $K^{\top, \bullet}_{T \times \Gm^{\rot}}(\pt)$; it is $2$-periodic and concentrated in even degrees; the affine cells give a basis of the degree zero part. The restriction maps induced by the skeletal inclusions $\sk_i \hookrightarrow \sk_{i'}$ are thus surjective; the inverse system in particular satisfies the Mittag--Leffler condition. This paving restricts to each affine Schubert variety $\Gr_{\leq \mu}$, where the same results hold. We deduce the following.
\begin{proposition}\label{proposition: completion of K-theory}
We have an identification of rings
 \begin{equation*}
  K^{\top, *}_{T \times \Gm^{\rot}}(\Gr_G) = \lim_{\mu \in X^+_{\bullet}} K^{\top, *}_{T \times \Gm^{\rot}}(\Gr_{\leq \mu}).   
\end{equation*} 
All terms are concentrated in the even part.
\end{proposition}
Since we do not loose anything by focusing on the degree zero piece $K^{\top, 0}_{T \times \Gm^{\rot}}(\Gr_G)$ of the $K$-theory ring, we do so from now on.


\subsubsection{GKM ring of affine Grassmannian.}\label{subsubsection: GKM ring of affine Grassmannian}
The action of the extended torus $T \times \Gm^{\rot}$ on $\Gr_G$ satisfies the GKM assumptions for integral topological $K$-theory; it thus may be described as an explicit subring of $K$-theory of the fixed points \cite{HHH05} and Proposition \ref{proposition: GKM of HHH}.
\begin{notation}\label{notation: GKM of Gr}
Let $\k$ be a base ring.
The GKM ring of the affine Grassmannian $\GKM^{0}_{T \times \Gm^{\rot}}(\Gr_G; \k) \subseteq \prod_{\lambda \in X_{\bullet}} \O(T \times \Gm^{\rot})$ is given by the subring
\begin{equation*}
    \GKM^{0}_{T \times \Gm^{\rot}}(\Gr_G; \k) 
    = \bigl\{ (f_{\lambda})_{\lambda} \in \prod_{\lambda \in X_{\bullet}} \O(T \times \Gm^{\rot}) \quad \big| \quad f_{\lambda} \equiv f_{s_{\beta + m\hbar} \cdot \lambda} \mod (1 - {\beta}q^m), \quad \forall \lambda, \forall s_{\beta + m\hbar} \bigr\},
\end{equation*}
where the congruence conditions run over all cocharacters $\lambda \in X_{\bullet}$ and all affine reflections $s_{\beta +m\hbar}$ associated to affine roots $\beta + m\hbar \in \Phi^{\bullet} \oplus \Z\hbar$.
\end{notation}

\begin{proposition}\label{proposition: GKM description of K-theory of Gr}
Over $\k=\Z$, we have
\begin{equation*}
K^{\top, 0}_{T \times \Gm^{\rot}}(\Gr_G; \k) = \GKM^{0}_{T \times \Gm^{\rot}}(\Gr_G; \k). 
\end{equation*}
\end{proposition}
\begin{proof}
The action of $T \times \Gm^{\rot} \acts \Gr_G$ satisfies the GKM assumptions for integral topological $K$-theory; this is the case for all Kac--Moody partial flag varieties by with respect to the natural torus \cite[\S 5]{HHH05}; let us sketch some details.

For $\Gr_G$, the fixed points of $T\times \Gm^{\rot}$ are given by $(t^{\lambda} \mid \lambda \in X_{\bullet})$ and we claim
\begin{equation*}
\begin{tikzcd}[row sep=5]
\GKM^{0}_{T \times \Gm^{\rot}}(\Gr_G; \k) \arrow[r, hookrightarrow] & \prod_{\lambda \in X_{\bullet}} \O(T \times \Gm^{\rot})\\
K^{\top, 0}_{T \times \Gm^{\rot}}(\Gr_G; \k) \arrow[u, equal] \arrow[r, hookrightarrow] & \prod_{\lambda \in X_{\bullet}} K^{\top, 0}_{T\times \Gm^{\rot}}(t^{\lambda}; \k) \arrow[u, equal].
\end{tikzcd}
\end{equation*}
Here, the Iwahori orbits on $\Gr_G$ provide a $T\times \Gm^{\rot}$-equivariant affine paving. Via the construction from Remark \ref{remark: GKM in the BB case}, the Assumptions \ref{assumption: HHH assumptions} are satisfied by \cite[\S 5, p.211]{HHH05}. The ring $K^{\top, 0}_{T \times \Gm^{\rot}}(\Gr_G; \k)$ thus has a GKM description by Proposition \ref{proposition: GKM of HHH}. Unraveling this gives precisely the ring from Notation \ref{notation: GKM of Gr} as in \cite{HHH05}. 
\end{proof}


\begin{remark}
Also see \cite{HHH05}, \cite[\S 2.2.3]{BBSV22}, \cite{Par17}, \cite[Appendix A]{Sit24}, \cite{Shi14} for the parallel overview of complexified cohomology of $\Gr_G$ in our situation.
\end{remark}





\begin{remark}
Given an affine Schubert variety $\Gr_{\leq \mu}$, the corresponding GKM ring is given by passing to the subset $\wt(\mu) \subseteq X_{\bullet}$ as
\begin{equation*}
    \GKM^{0}_{T \times \Gm^{\rot}}(\Gr_{\leq \mu}; \k) 
    = \bigl\{ (f_{\lambda}) \in \prod_{\lambda \in \wt(\mu)} \O(T \times \Gm^{\rot}) \quad \big| \quad f_{\lambda} \equiv f_{\lambda'} \mod (1 - {\beta}q^m) \bigr\},
\end{equation*}
where the conditions run over all pairs $\lambda, \lambda' \in \wt(\mu)$ such that $\lambda' = s_{\beta + m\hbar} \cdot \lambda$ for a simple reflection associated to some affine root $\beta + m\hbar \in \Phi^{\bullet} \oplus \Z \hbar$. We again have
\begin{equation*}
K^{\top, 0}_{T \times \Gm^{\rot}}(\Gr_{\leq \mu}; \k) = \GKM^{0}_{T \times \Gm^{\rot}}(\Gr_{\leq \mu}; \k). 
\end{equation*}
The restriction maps from $K^{\top, 0}_{T \times \Gm^{\rot}}(\Gr_G; \k)$ on the left correspond to restriction along $\wt(\mu) \subseteq X_{\bullet}$ on the right. Note that $\Gr_{\leq \mu}$ are GKM spaces in the more standard sense -- the number of fixed points is finite.
\end{remark}


Our aim is to describe this $K$-theory ring, in other words explicitly solve the system of congruences defining $\GKM^0_{T \times \Gm^{\rot}}(\Gr; \k)$ when $\k=\C$. However, it is nontrivial to do so directly, and we apply other topological tricks. In any case, our Theorem \ref{theorem: equivariant K-theory as deformation to normal cone -- G sssc} below may be regarded as a solution to this purely combinatorial problem. 







\subsection{Quasi-diagonals at roots of unity}\label{section: quasi-diagonals at roots of unity}
Let $\k$ be any coefficient ring and $G$ a split reductive group over $\k$ with maximal torus $T$ and Weyl group $W$.
In \S \ref{subsubsection: relationship to fibers of Gamma}, \S \ref{subsubsection: generatirs of the ideal sheaf of Zl} we assume $G=G^{\sc}$ to be simply connected. 

\subsubsection{Power $\ell$ map.}
Consider $T \times T$; we denote respectively $t_1, \dots, t_n$ and $s_1, \dots, s_n$ coordinates on the first and second copy. 
%
Denote
\begin{equation*}
    (-)^{\ell} : T \to T, \qquad t \mapsto t^{\ell}
\end{equation*}
the $\ell$-th power map, coming from the group structure on $T$. In coordinates, it sends $t_i \mapsto t_i^{\ell}$. Since the $W$-action on $T$ is by group automorphisms, the $\ell$-th power map is $W$-equivariant. It thus induces a map
\begin{equation*}
    (-)^{\ell} : T \sslash W \to T \sslash W.
\end{equation*}
The corresponding map $\O(T \sslash W) \to \O(T \sslash W)$ on rings of functions is induced by the substitution $f(t_1, \dots, t_n) \mapsto f(t^{\ell}_1, \dots, t^{\ell}_n)$. Similarly for the quotient by any subgroup of $W$.

\subsubsection{Quasi-diagonals at roots of unity.}
For any $\ell \in \N$, we define the corresponding \textit{quasi-diagonal} by the pullback
\begin{equation*}
\begin{tikzcd}
Z^{T \sslash W, T\sslash W}_{\ell} \arrow[d] \arrow[r]  & T \sslash W \arrow[d, "(-)^{\ell}"]  \\
T \sslash W \arrow[r, "(-)^{\ell}"] & T \sslash W
\end{tikzcd}
\end{equation*}
It has trivial derived structure by flatness of $(-)^{\ell}$.
It is a closed subscheme of $T \sslash W \times T \sslash W$ via the pullback square
\begin{equation*}
    \begin{tikzcd}
        Z^{T \sslash W, T\sslash W}_{\ell} \arrow[r, hook] \arrow[d] & T\sslash W \times T\sslash W \arrow[d, "(-)^{\ell} \times (-)^{\ell}"] \\
        T\sslash W \arrow[r, hook, "\Delta"] & T\sslash W \times T\sslash W
    \end{tikzcd}
\end{equation*}
On its functor of points, it parametrizes
\begin{equation*}
Z^{T \sslash W, T\sslash W}_{\ell}: R \mapsto \{ (x, y) \in (T\sslash W \times T\sslash W)(R) \mid x^{\ell} = y^{\ell} \}.
\end{equation*}

\subsubsection{Variants.}
More generally, for any two subgroups of $ W_1, W_2 \leq W$, we define the variant $Z^{W_1, W_2}_{\ell} = Z_{\ell}^{T \sslash W_1, T\sslash W_2}$ by either of the pullback squares
\begin{equation*}
\begin{tikzcd}
Z_{\ell}^{T \sslash W_1, T\sslash W_2} \arrow[d] \arrow[r]  & T \sslash W_2 \arrow[d, "(-)^{\ell}"]  \\
T \sslash W_1 \arrow[r, "(-)^{\ell}"] & T \sslash W
\end{tikzcd}
\qquad
\begin{tikzcd}
Z^{T \sslash W_1, T\sslash W_2}_{\ell} \arrow[r, hook] \arrow[d] & T\sslash W_1 \times T\sslash W_2 \arrow[d, "(-)^{\ell} \times (-)^{\ell}"] \\
T\sslash W \arrow[r, hook, "\Delta"] & T\sslash W \times T\sslash W.
\end{tikzcd}
\end{equation*}
It is a closed subscheme of $T\sslash W_1 \times T\sslash W_2$, obtained by the base change of $Z^{T \sslash W, T \sslash W}_{\ell}$ along the faithfully flat map $T\sslash W_1 \times T\sslash W_2 \to T\sslash W \times T\sslash W$.
%
On its functor of points, it parametrizes
\begin{equation*}
Z^{T\sslash W_1 \times T\sslash W_2}_{\ell}(R) = \{ (x, y) \in (T\sslash W_1 \times T\sslash W_2)(R) \mid x^{\ell} W = y^{\ell} W \}.    
\end{equation*}

\begin{notation}\label{notation: quasi-diagonals}
In particular, we will use the following shorter notation for two special cases:
\begin{align*}
Z'_{\ell} & := Z^{T, T}_{\ell} \hookrightarrow T \times T,   \\
Z_{\ell} & := Z^{T, T\sslash W}_{\ell} \hookrightarrow T \times T \sslash W.
\end{align*}
\end{notation}

\begin{notation}
Furthermore, given $\zeta \in \mu_{\infty}$ of primitive order $\ell$, we will also use the notation $Z_{\zeta} = Z_{\ell}\hookrightarrow T \times T \sslash W$. Given any point $x \in T$ in the first copy of the torus, we denote its fiber over $x$ by $Z_{(x, \zeta)} = Z_{(x, \ell)} \hookrightarrow T \sslash W$.    
\end{notation}




\subsubsection{Relationship to fibers of $\Gamma$.}\label{subsubsection: relationship to fibers of Gamma}
Assume now $\k$ is an algebraically closed field of characteristic zero. Assume $G$ is simply connected.
Given $\zeta \in \mu_{\infty}$ of primitive order $\ell$, we have just defined the subscheme $Z_{\zeta}$ and its fibers $Z_{(x, \zeta)}$.
%
These are fully described by specializations of the sub-ind scheme $\overline{\Gamma} \hookrightarrow T \times \Gm^{\rot} \times T \sslash W$ from \S \ref{section: cocharacter graphs}.


\begin{lemma}\label{lemma: Z is Gamma}
Let $(x, \zeta) \in T \times \Gm^{\rot}$. Suppose $\zeta$ is a root of unity of primitive order $\ell$. On the underlying points, we have
\begin{equation*}
  |Z'_{\ell}| = |{\Gamma}_{\zeta}|,  \qquad \qquad |Z'_{(x, \ell)}| = |{\Gamma}_{(x, \zeta)}|, \qquad \qquad
  |Z_{\ell}| = |\overline{\Gamma}_{\zeta}|,  \qquad \qquad |Z_{(x, \ell)}| = |\overline{\Gamma}_{(x, \zeta)}|.
\end{equation*}
\end{lemma}
\begin{proof}
   It is enough to treat the $Z'_{\ell}$-variants, since the $Z_{\ell}$-variant follows by taking quatient by $W$. Moreover, it is enough to focus on the case of $Z'_{(x, \ell)}$ (which is a fiberwise refinement of $Z'_{\ell}$). Over any $(x, \zeta)$, both sets have a transitive action of $W_{\aff}$ and base point with the same stabilizer.
\end{proof}

\begin{remark}
    More explicitly, we may just compute inside $T \times T$ as follows
    \begin{align*}
       |Z'_{\ell}|  
       &=  \{ (x, y) \mid \exists w \in W: x^{\ell} = w(y^{\ell}) \} \\ 
       &= \{ (x, y) \mid \exists w \in W: (x \cdot w(y)^{-1})^{\ell} = 1 \} \\
       &= \{ (x, y) \mid \exists \lambda \in X_{\bullet}, \exists w \in W: x \cdot w(y)^{-1} = \lambda(\zeta) \} \\
       &= |\Gamma_{\zeta}| 
    \end{align*}
    The case of $Z'_{(x, \ell)}$ is the same computation with fixed $x$.
\end{remark}

\begin{remark}
The statement of Lemma \ref{lemma: Z is Gamma} holds on the level of reduced schemes: it identifies $Z_{\ell} = \overline{\Gamma}_{\zeta}^{\red}$ and $Z_{(x, \ell)} = \overline{\Gamma}_{(x, \zeta)}^{\red}$. 

However, while $Z_{\ell}$ is reduced by definition, the fiber $\overline{\Gamma}_{\zeta}$ will acquire some ind-unreduced structure from the $q$-direction. In fact, one may heuristically think about the normal cone to $Z_{\ell}$ as a counterpart of this ind-unreduced structure.
\end{remark}





\subsubsection{Generators of the ideal sheaf of $Z_{\ell}$.}\label{subsubsection: generatirs of the ideal sheaf of Zl}
Let again $\k$ be a general coefficient ring.
Consider the ring $\O(T \sslash W)$ over $\k$. It naturally identifies with the representation ring $R(G; \k)$. Denote $e_1, \dots, e_n \in \O(T \sslash W)$ the elements corresponding to the fundamental representations of $G$. Since $G$ is simply connected or $GL_n$, these freely generate $\O(T \sslash W)$ by \S \ref{subsubsection: coordinates and invariant polynomials}.

\begin{lemma}\label{lemma: explicit generators for the ideal sheaf of Zl}
The ideal sheaf $\eI_{Z_{\ell}}$ of the closed subscheme $Z^{T\sslash W \times T \sslash W}_{\ell} \hookrightarrow T\sslash W \times T \sslash W$ is generated by
\begin{equation*}
 f_i = e_i(t_1^{\ell}, \dots, t_n^{\ell}) - e_i(s_1^{\ell}, \dots, s_n^{\ell}), \qquad i = 1, \dots, n    
\end{equation*}
The same works for the ideal sheaf $\eI_{Z_{\ell}}^{T \sslash W_1, T \sslash W_2}$ of $Z_{\ell}^{T \sslash W_1, T \sslash W_2} \hookrightarrow T \sslash W_1 \times T \sslash W_2$ for any $W_1, W_2 \leq W$.
\end{lemma}
\begin{proof}  

By definition, $Z_{\ell}$ fits into the pullback square
\begin{equation}\label{diagram: fiber square for Zl, variant II}
    \begin{tikzcd}
        Z_{\ell} \arrow[r, hook] \arrow[d] & T\sslash W \times T\sslash W \arrow[d, "(-)^{\ell} \times (-)^{\ell}"] \\
        T\sslash W \arrow[r, hook, "\Delta"] & T\sslash W \times T\sslash W
    \end{tikzcd}
\end{equation}
Note that $T\sslash W$ is smooth (in fact an affine space) and the map $(-)^{\ell}: T\sslash W \to T \sslash W$ is finite (which is obvious before taking the quotient). The same is true for $(-)^{\ell} \times (-)^{\ell}: T\sslash W \times T\sslash W \to T\sslash W \times T\sslash W$, which is thus flat by miracle flatness \cite[Tag 00R4]{Sta}. Since it is surjective, it is in fact faithfully flat.

The closed immersions given by the horizontal arrows in \eqref{diagram: fiber square for Zl, variant II} induce short exact sequence
\begin{equation*}
\begin{tikzcd}
0 \arrow[r] & \eI_{Z_{\ell}} \arrow[r] & \O_{T\sslash W \times T\sslash W} \arrow[r] & \O_{Z_{\ell}} \arrow[r] & 0 \\
0 \arrow[r] & \eI_{T\sslash W} \arrow[r] \arrow[u] & \O_{T\sslash W \times T\sslash W} \arrow[r] \arrow[u] & \O_{T\sslash W} \arrow[u] \arrow[r] & 0 
\end{tikzcd}
\end{equation*}
The upper sequence is given by tensoring the bottom one along $(-)^{\ell} \times (-)^{\ell}: \O_{T\sslash W \times T\sslash W} \hookrightarrow \O_{T\sslash W \times T\sslash W}$. Indeed, this is true for the right-hand term by \eqref{diagram: fiber square for Zl, variant II} and tautological for the middle term, hence true for the left-hand term by flatness. In coordinates of $T \times T$, this map is given by $t_i \mapsto t_i^{\ell}$ and $s_i \mapsto s_i^{\ell}$ for any $i = 1, \dots, n$. 

Now, the ideal sheaf $\eI_{T\sslash W}$ of the diagonal in $T\sslash W \times T\sslash W$ is generated by the elements $f_j = e_j(t_1, \dots, t_n) - e_j(s_1, \dots, s_n)$ for $j = 1, \dots, n$ (as these are affine spaces with coordinates $e_j$.) Therefore, the ideal sheaf $\eI_{Z_{\ell}}$ is generated by their images
\begin{equation*}
    \eI_{Z_{\ell}} = \langle e_j(t^{\ell}_1, \dots, t^{\ell}_n) - e_j(s^{\ell}_1, \dots, s^{\ell}_n) \mid j = 1, \dots, n \rangle.
\end{equation*}
as an ideal inside $\O_{T\sslash W \times T \sslash W}$. 


The general case follows, since $Z_{\ell}^{T \sslash W_1, T \sslash W_2} \hookrightarrow T \sslash W_1 \times T \sslash W_2$ is given by the base change of $Z_{\ell}^{T \sslash W, T \sslash W} \hookrightarrow T \sslash W \times T \sslash W$ along the faithfully flat map $T \sslash W_1 \times T \sslash W_2 \to T \sslash W \times T \sslash W$.
\end{proof}



\subsection{Deformations of quasi-diagonals to normal cones}\label{section: deformations of quasi-diagonals to normal cones}

\subsubsection{Deforming to the normal cone at roots of unity.}
Assume $G=G^{\sc}$ is simply connected or $GL_n$.
Recall the general notation for deformations to normal cones and affine blowups from \S \ref{section: deformations to normal cones and affine blowups}. To simplify the manipulation of roots of unity, assume $\k$ is an algebraically closed field of $\chara \k = 0$.

\begin{definition}\label{definition: deformation to normal cone at roots of unity}
Take the scheme $T \times T \sslash W \times \Gm^{\rot}$, viewed as a family over $\Gm^{\rot} \hookrightarrow \A^1$. Consider the family of subschemes $\eZ = (Z_{\zeta} \mid \zeta \in \mu_{\infty})$ of its fibers over varying $\zeta \in \mu_{\infty} \hookrightarrow \Gm^{\rot}$. We consider the affine blowup of this family
\begin{equation*}
    \eN_{\mu_{\infty}}(\eZ, T \times T \sslash W) := \Bl^{\circ}_{\eZ}(T \times T \sslash W \times \Gm^{\rot}).
\end{equation*}
\end{definition}

In terms of deformations to normal cones, we can restate the above as follows. Start with the scheme $T \times T\sslash W$. Then $\eN_{\mu_{\infty}}(\eZ, T \times T \sslash W)$ is the family over $\Gm^{\rot}$ which is the constant family with fiber $T \times T\sslash W$ away from roots of unity and the deformation to the normal cone of $Z_{\ell} \hookrightarrow T \times T\sslash W$ on each (formal) neighborhood of each root of unity $\zeta$ of primitive order $\ell \in \N$. In other words, we have pullback squares
\begin{equation*}
    \begin{tikzcd}
      C(Z_{\zeta}, T \times T \sslash W) \arrow[d] \arrow[r]&  \eN_{\mu_{\infty}}(\eZ, T \times T \sslash W) \arrow[d] & \arrow[l] \arrow[d] T \times T \sslash W \times (\Gm^{\rot} - \mu_{\infty}) \\
      \zeta \arrow[r] & \Gm^{\rot} & \arrow[l] \Gm^{\rot} - \mu_{\infty}. 
    \end{tikzcd}
\end{equation*}
Our family $\eN_{\mu_{\infty}}(Z, T \times T \sslash W)$ is given the usual Rees constructions for the filtered algebras $(\O_{T \times T \sslash W}, \eI^{\bullet}_{Z_{\ell}})$ at each $\zeta$ glued together along $\Gm^{\rot} - \mu_{\infty}$, where all of them naturally trivialize.





\subsubsection{Explicit description of the deformation.}
Since $T \times T \sslash W$ is affine, the scheme $\eN_{\mu_{\infty}}(\eZ, T \times T \sslash W)$ is affine with ring of functions
\begin{equation*}
\O(\eN_{\mu_{\infty}}(\eZ, T \times T \sslash W))
=
\O(T \times (T \sslash W))[q^{\pm 1}][\tfrac{f}{q - \zeta} \mid (\zeta, f) \in \Xi]
\end{equation*} 
where $\Xi$ is the set of all pairs $(\zeta, f) \in \k^{\times} \times \O(T  \times (T \sslash W))$ with $\zeta^{\ell} = 1$ and $f \in \eI_{Z_{\ell}}$ for a common $\ell \in \N$.
Note that this algebra is \textit{not} finitely generated over $\k$.    
Alternatively, grouping all $\ell$-torsion roots of unity together, we can rewrite it as follows.
\begin{lemma}\label{claim: putting ell-th roots of unity together} We have
\begin{equation*}
\O(\eN_{\mu_{\infty}}(\eZ, T \times T \sslash W))
=
\O(T \times (T \sslash W))[q^{\pm 1}][\tfrac{f}{1 - q^{\ell}} \mid \ell \in \N, \ f \in \eI_{Z_{\ell}}].
\end{equation*}
\end{lemma}
\begin{proof}
Recall $\k = \overline{\k}$ of $\chara \k = 0$.
Write $(q^{\ell}-1) = (q-\zeta_0)\cdots (q-\zeta_{\ell-1})$, $\zeta_i \in \mu_{\ell}$. Clearly, each
$\tfrac{f}{\zeta -q} = \tfrac{f}{1 - q^{\ell}} \cdot \prod_{i'\neq i} (q-\zeta_{i'})$ is generated by $\tfrac{f}{1 - q^{\ell}}$.
%
On the other hand, all the monomials $(q-\zeta_i)$ are pairwise coprime in $\k[q^{\pm 1}]$. By the Chinese remainder theorem, we may express
$1 = \sum_{i=0}^{\ell-1} h_i \cdot \prod_{i'\neq i} (q-\zeta_{i'})$ for some $h_i \in \k[q^{\pm 1}]$, 
hence
$1 = \sum_{i=0}^{\ell-1} h_i \cdot (q^{\ell} - 1) \tfrac{1}{q - \zeta_i}.$
Consequently, $\tfrac{f}{1 - q^{\ell}} = \sum_{i=0}^{\ell-1} h_i (q^{\ell} - 1) \tfrac{f}{\zeta - q}$ is generated by the elements $\tfrac{f}{\zeta - q}$.
\end{proof}



\begin{remark}
For any $\ell \in \N$, we denote the corresponding cyclotomic polynomial by
$\sigma_{\ell}(q) \in \Z[q^{\pm 1}]$.    
Since we work over an algebraically closed field $\k$ of $\chara \k = 0$, we also have
\begin{equation*}
  \O(T \times (T \sslash W))[q^{\pm 1}][\tfrac{f}{\sigma_{\ell}(q)} \mid \ell \in \N, \ f \in \eI_{Z_{\ell}}]  
\end{equation*}
by the same token as the lemma above.
\end{remark}





\subsubsection{Fundamental Bezrukavnikov classes.}
To generate the functions on the deformation, the following distinguished elements suffice. Recall that we are assuming $G$ simply connected or $GL_n$.
\begin{definition}
We define the \textit{fundamental Bezrukavnikov classes} as
\begin{equation*}
    b_{j, \ell} := \tfrac{f_j}{1-q^{\ell}} := \tfrac{e_j(t_1^{\ell}, \dots, t_n^{\ell}) - e_j(s_1^{\ell}, \dots, s_n^{\ell})}{1 - q^{\ell}}, \qquad \forall j \in \{ 1, \dots, n\}, \forall \ell \in \N. 
\end{equation*}
\end{definition}

\begin{lemma}\label{lemma: fundamental bezrukavnikov classes}
To generate $\O(\eN_{\mu_{\infty}}(\eZ, T\sslash W \times T \sslash W))$ as an algebra over $\O(T\sslash W \times \Gm^{\rot} \times T \sslash W)$, it suffices to add the infinite set of ring generators
\begin{equation*}
    \left( b_{j, \ell} \mid \forall j \in \{ 1, \dots, n\}, \forall \ell \in \N \right). 
\end{equation*}
\end{lemma}
\begin{proof}
By definition, for each $\zeta \in \mu_{\infty}$ and every $f \in \eI_{Z_{\zeta}}$, we need to add the generator $\tfrac{f}{q-\zeta}$. For fixed $\zeta$, it is enough to only consider $b_{j, \zeta} := \tfrac{f_j}{q-\zeta}$ for $f_j = e_j(t_1^{\ell}, \dots, t_n^{\ell}) - e_j(s_1^{\ell}, \dots, s_n^{\ell})$, because $f_1, \dots, f_j$ generate $\eI_{Z_{\ell}}$ by Lemma \ref{lemma: explicit generators for the ideal sheaf of Zl}.

By the same coprimality argument as in Lemma \ref{claim: putting ell-th roots of unity together}, we may group all the $\ell$-torsion roots of unity together -- the elements $b_{j, \zeta}$ and $b_{j, \ell}$ generate the same ring.
\end{proof}

\begin{remark}
Consider the base changed variant $Z^{T\sslash W_1, T \sslash W_2}_{\ell} \hookrightarrow T\sslash W_1 \times T \sslash W_2$. The same set $(b_{j, \ell})$ then generates $\O(N_{q=\zeta}(Z^{T\sslash W_1, T \sslash W_2}_{\ell}, T\sslash W_1 \times T \sslash W_2))$ over $\O(T\sslash W_1 \times \Gm^{\rot} \times T \sslash W_2)$: this follows by base change along the flat map $T\sslash W_1 \times T \sslash W_2 \to T\sslash W \times T \sslash W$.  
\end{remark}



\subsubsection{Completions of the deformation over a root of unity.}

The following lemma is the algebraic analogue of a local comparison between equivariant $K$-theory and cohomology of fixed points.

\begin{lemma}\label{lemma: completions of the deformation}
Let $(x, \zeta) \in T \times \Gm^{\rot}$ with $\zeta \in \mu_{\infty}$. Then
\begin{equation*}
    N_{q=\zeta}(Z_{\ell}, T \times T\sslash W))^{\wedge}_{(x, \zeta)} = \coprod_{\vartheta \in |Z_{(x, \zeta)}|} N_{h=0}(\t \times_{\t \sslash W_{[x, \zeta]}} \t \sslash W_{\vartheta}, \t \times \t \sslash W_{\vartheta})^{\wedge}_{(0, 0)}
\end{equation*}
\end{lemma}
\begin{proof}
We first consider the case of $Z'_{\zeta} \hookrightarrow T \times T$. Recall we have already identified $|Z'_{\zeta}|=|\overline{\Gamma}_{\zeta}|$ and $|Z'_{(x, \zeta)}| = |\Gamma_{(\zeta, x)}| $  in Lemma \ref{lemma: Z is Gamma}.

Decompose $Z_{(x,\zeta)}'$ into its connected components ${Z'}_{(x, \zeta)}^{\vartheta}$ labeled by $\vartheta \in |Z_{(x, \ell)}'|$. We make the following identifications:
\begin{align*}
  N_{q=\zeta}(Z'_{\zeta}, T \times T))^{\wedge}_{(x, \zeta)}
  &= \coprod_{\vartheta \in |Z'_{(x, \zeta)}|} N_{q=\zeta}({Z'}_{(x, \zeta)}^{\vartheta}, T \times T)^{\wedge}_{(x, \zeta)} \\
  &= \coprod_{\vartheta \in |Z'_{(x, \zeta)}|} N_{q=1}(T \times_{T \sslash W_{[x, \zeta]}} T, T \times T)^{\wedge}_{(1, 1)} \\
  &= \coprod_{\vartheta \in |Z'_{(x, \zeta)}|} N_{h=0}(\t \times_{\t \sslash W_{[x, \zeta]}} \t, \t \times \t)^{\wedge}_{(0, 0)}.
\end{align*}

The first line holds since the deformation to the normal cone is Zariski local and compatible with connected components, so we only need to study each of ${Z'}^{\vartheta}_{(x, \zeta)} \hookrightarrow T \times T$ separately. 

Now, the irreducible components of $|Z'_{\ell}|$ are labeled by $W_{\aff} / \ell X_{\bullet}$. Each of these connected components is a copy of $T$. On the other hand, the irreducible components of its fiber $|Z'_{(x, \zeta)}|$ are labeled by the set $W_{\aff} / \Stab_{W_{\aff}}{(x, \zeta, x)}$. By \eqref{equation: stabilizer exact sequences for zeta root of unity}, Remark \ref{remark: colisions of components of Gamma}, the collisions of the connected components of $Z'_{\ell}$ over $x$ are described by the $W_{[x, \zeta]}$-torsor of finite sets
\begin{equation*}
  W_{\aff} / \ell X_{\bullet} \twoheadrightarrow W_{\aff} / (\ell X_{\bullet} \rtimes W_{[x, \zeta]}, \bullet_{\omega}  )
\end{equation*}
where $W_{[x, \zeta]}$ acts by the $\omega$-shifted right action.
The colliding copies of $T$ at $\vartheta$ are labeled by $W_{[x, \zeta]}$, giving the second identification.

The third line is the standard identification of formal neighborhood of the identity in a reductive group with its Lie algebra (also see \cite[(2.17)]{BBSV22}). This concludes the discussion for $Z'_{\ell}$.


Now, consider the action of $W$ on the left hand side, induced from the tautological action on the second copy of $T$. It translates to the right-hand side as follows. Firstly, it permutes the connected components: given a component labeled by $\vartheta \in |Z'_{(x, \zeta)}|$, its orbit under $W \acts |Z'_{(x, \zeta)}|$ identifies with $W/W_{\vartheta}$. This cuts down the labeling set to $|Z_{(x, \zeta)}|$ as in Lemma \ref{lemma: fiber of cocharacter graphs as lattice quotient}. Furthermore, fixing any such $\vartheta$, we are left with the natural action of $W_{\vartheta}$ on the closed immersion ${Z'}_{(x, \zeta)}^{\vartheta} \hookrightarrow T \times T$.

Altogether, taking the quotient by $W$ on both sides gives the desired equality
\begin{equation*}
    N_{q=\zeta}(Z_{\ell}, T \times T\sslash W))^{\wedge}_{(x, \zeta)} = \coprod_{\vartheta \in |Z_{(x, \zeta)}|} N_{h=0}(\t \times_{\t \sslash W_{[x, \zeta]}} \t \sslash W_{\vartheta}, \t \times \t \sslash W_{\vartheta})^{\wedge}_{(0, 0)}.
\end{equation*}
\end{proof}

\begin{remark}
More generally, the final step of the above prove shows that for $(x, \zeta) \in T \times \Gm^{\rot}$ with $\zeta \in \mu_{\infty}$ and any pair of subgroups $W_1, W_2 \leq W$ we have
\begin{equation*}
    N_{q=\zeta}(Z^{W_1, W_2}_{\ell}, T\sslash W_1 \times T\sslash W_2))^{\wedge}_{(x, \zeta)} = \coprod_{\vartheta \in |Z^{W_1, W_2}_{(x, \zeta)}|} N_{h=0}(\t \sslash W_1 \times_{\t \sslash W_{[x, \zeta]}} \t \sslash W_{2, \vartheta}, \t \sslash W_1 \times \t \sslash W_{2, \vartheta})^{\wedge}_{(0, 0)}.
\end{equation*}
\end{remark}


\subsubsection{Strict transforms of cocharacter graphs.}
Finally, it is quite useful to keep track of how the cocharacter graphs transform in the affine blowup.



\begin{notation}
Denote 
\begin{equation*}
 \overline{\Upsilon} := \overline{\Gamma}^{\st} \hookrightarrow \eN_{\mu_{\infty}}(\eZ, T \times T\sslash W) = \Bl^{\circ}_{\eZ}(T \times T\sslash W \times \Gm^{\rot})   
\end{equation*}
the strict transform of $\overline{\Gamma} \hookrightarrow T \times \Gm^{\rot} \times T \sslash W$.
This is a closed sub-ind-schemes by \S \ref{section: embedding of strict transform into affine blowup}.
%
The same construction works with other variants -- in particular, we denote ${\Upsilon} := {\Gamma}^{\st} \hookrightarrow \eN_{\mu_{\infty}}(\eZ', T \times T) = \Bl^{\circ}_{\eZ'}(T \times T \times \Gm^{\rot})$ the strict transform of $\Gamma \hookrightarrow T \times \Gm^{\rot} \times T$.
\end{notation}
%
We have the following situation
\begin{equation}\label{equation: Gamma and Upsilon}
\begin{tikzcd}
\Bl^{\circ}_{\eZ}(T \times \Gm^{\rot} \times T \sslash W) \arrow[d] \arrow[r, hookleftarrow] & \overline{\Upsilon}  \arrow[r, hookleftarrow] \arrow[d] & \overline{\Upsilon}_{\wt(\mu)} \arrow[d] \\
T \times \Gm^{\rot} \times T \sslash W \arrow[r, hookleftarrow] & \overline{\Gamma} \arrow[r, hookleftarrow] & \overline{\Gamma}_{\wt(\mu)}.   
\end{tikzcd}
\end{equation}
The horizontal maps are closed immersions of ind-schemes, the outer terms are actual schemes.





By \S \ref{section: strict transforms}, each fiber $\overline{\Upsilon}$ is given by the blowup of $\overline{\Gamma}$ at the ind-subscheme $\eZ$. In particular, restricting to a formal neighborhood of $(x, \zeta)$, we see that
\begin{equation*}
    \overline{\Upsilon}^{\wedge}_{(x, \zeta)} = \Bl_{Z^{\wedge}_{\zeta}} \overline{\Gamma}^{\wedge}_{(x, \zeta)}.
\end{equation*}


\begin{remark}\label{remark: Upsilon and GKM}
One can trace where $\overline{\Upsilon}$ and $\Upsilon$ go under the identification of Lemma \ref{lemma: completions of the deformation}. Firstly, it decomposes into connected components $\overline{\Upsilon}^{\wedge}_{(x, \zeta)} = \coprod_{\vartheta \in |Z_{(x, \zeta)}|} \overline{\Upsilon}^{\vartheta, \wedge}_{(x, \zeta)}$. Each of the closed immersions
\begin{equation*}
\overline{\Upsilon}^{\vartheta, \wedge}_{(x, \zeta)} \hookrightarrow N_{h=0}(\t \times_{\t \sslash W_{[x, \zeta]}} \t \sslash W_{\vartheta}, \t \times \t \sslash W_{\vartheta})^{\wedge}_{(x, \zeta)}  
\end{equation*}
then identifies as a hyperplane arrangement of copies of $(\t\times \Ga)^{\wedge}_{(0, 0)}$ inside $(\t \times \Ga \times \t \sslash W_{\vartheta})^{\wedge}_{(0, 0)}$ intersecting in the origin, labeled by those elements $\lambda \in X_{\bullet}$ which lie in the corresponding coset $\vartheta \in X_{\bullet}/\Stab_{W_{\aff}}{(x, \zeta, x)}$, see also Lemma \ref{lemma: fiber of cocharacter graphs as lattice quotient}. Depending on the genericity of $\zeta$, the precise arrangement is that of Remarks \ref{remark: GKM for cohomology of finite flag varieties} and \ref{remark: informal GKM localization}, applied to the ambient reductive group $L_{[x, \zeta]}$ with the parabolic subgroup $P_{\vartheta}$.
\end{remark}



\subsubsection{Modifications in the reductive case.}\label{subsubsection: modifications in the reductive case}
While the subschemes $Z_{\zeta} = Z_{\ell} \hookrightarrow T \times T\sslash W$ from \S \ref{section: quasi-diagonals at roots of unity} make sense for any reductive $G$, they have to be modified for our purposes. For start, Lemma \ref{lemma: explicit generators for the ideal sheaf of Zl} is specific to $T=T^{\sc}$. Related modification is needed to generalize Lemma \ref{lemma: Z is Gamma} to the setup from \S \ref{section: cocharacter graphs for reductive groups}.

In view of this, we extend Definition \ref{definition: deformation to normal cone at roots of unity} of $\eZ$ and $\Bl^{\circ}_{\eZ}(T \times \Gm^{\rot} \times T\sslash W)$ to the case of reductive $G$ as follows. 
We consider the family $\widetilde{\eZ}$ of subschemes $(T \times_{T^{\ad}} (Z^{\sc}_{\zeta}/\pi_1(G^{\ad}))) \times \pi_1(G)) \subseteq T \times T^{\sc} \sslash W \times \pi_1(G)$. In words, we take one copy of $\eZ$ for $G^{\sc}$ in each connected component, descend it along the quotient by $\pi_1(G^{\ad})$, and finally pull back to the given torus $T$ as in \S \ref{section: action of the fundamental group}, Remarks \ref{remark: identification for Gamma in semisimple case} and \ref{remark: identification for Gamma in general reductive case}.

With this definition, the affine blowup $\Bl^{\circ}_{\widetilde{\eZ}}(T \times \Gm^{\rot} \times T^{\sc} \sslash W \times \pi_1(G))$ is centered at the family of reduced subschemes underlying $\Gamma^{T, T^{\sc}}_{\zeta}$ for $\zeta \in \mu_{\infty}$ from \S \ref{section: action of the fundamental group}. Similarly for actions of other $W_1, W_2 \leq W$ on the two copies of $T$.





\subsection{A combinatorial map}\label{section: map to equivariant K-theory: combinatorics}
We construct a comparison map from the completed ring of functions on the above deformation to the GKM ring of the affine Grassmannian. In this section, we work purely combinatorially. The parallel topological meaning is explained in \S \ref{section: map to equivariant K-theory: topology}. 

Let $\k$ be a coefficient ring. To simplify the manipulation of roots of unity in our deformation, we assume $\k$ is an algebraically closed field of $\chara \k = 0$. While the discussion makes sense for split reductive $G$, we focus on the relevant case of $G=G^{\sc}$ simply connected or $GL_n$.


\subsubsection{Components of the localization map.}
We define the localization map $$\loc: \O(T \times \Gm^{\rot} \times T \sslash W) \xrightarrow{\loc} \prod_{\lambda \in X_{\bullet}} \O(T \times \Gm^{\rot})$$ componentwise as
\begin{align}\label{equation: combinatorial localization map}
   \loc_{\lambda}: \O(T \times \Gm^{\rot} \times T \sslash W) &\xrightarrow{\loc} \prod_{\lambda \in X_{\bullet}} \O(T \times \Gm^{\rot}) \xrightarrow{\pr_{\lambda}} \O(T \times \Gm^{\rot}), \\   
   \gamma_{\lambda}: f(x, \zeta, y) &\mapsto f(x, \zeta, \lambda(\zeta) \cdot x).
\end{align}
%
Consider the closed embedding $\overline{\Gamma}_{\lambda} \hookrightarrow T \times \Gm^{\rot} \times (T \sslash W)$ and identify $\overline{\Gamma}_{\lambda} = T \times \Gm^{\rot}$ via the first two factors. The components $\loc_{\lambda}$ are then geometrically given by restriction to the irreducible components $\overline{\Gamma}_{\lambda}$ of $\overline{\Gamma}$, namely
\begin{equation}
   \loc_{\lambda}: \O(T \times \Gm^{\rot} \times T \sslash W) \twoheadrightarrow \O(\overline{\Gamma}_{\lambda}).
\end{equation}

We can now combinatorially check that the image of $\gamma$ satisfies the GKM conditions.
\begin{lemma}\label{lemma: GKM localization map combinatorial}
The image of the map $\loc$ from \eqref{equation: combinatorial localization map} lands in the subring $\GKM^{0}_{T \times \Gm^{\rot}}(\Gr_G; \k)$.     
\end{lemma}
\begin{proof}
We directly check the GKM conditions. To simplify notation, observe that each function from the base $\O(T \times \Gm^{\rot})$ restrict the same to each $\overline{\Gamma}_{\lambda}$, hence satisfies the GKM conditions trivially. It is thus enough to check on functions $f \in \O(T\sslash W)$ from the second copy. We identify $f$ with the corresponding $W$-invariant function on $T$ via the embedding $\O(T\sslash W) \hookrightarrow \O(T)$; the restriction to $\overline{\Gamma}_{\lambda} \hookrightarrow T\sslash W$ now corresponds to restriction to $\Gamma_{\lambda, 1} \hookrightarrow T$.

Consider any affine simple reflection $s_{\beta + m\hbar}$. 
Let $\lambda, \lambda' \in X_{\bullet}$ be such that $\lambda' = s_{\beta+ m\hbar} (\lambda)$, or equivalently
\begin{equation}\label{equation: combinatorial verification eq 0}
\lambda' = s_{\beta}(\lambda) \cdot (\widecheck{\beta})^m.   
\end{equation}
Assume that $(x, \zeta) \in T \times \Gm^{\rot}$ is in the kernel of the character $\beta+m\hbar$. In other words, this means 
\begin{equation}\label{equation: combinatorial verification eq 1}
\beta(x) \cdot \zeta^m = 1,  
\end{equation}
or equivalently $\beta(x) = \zeta^{-m}$.
Consequently, we have 
\begin{equation}\label{equation: combinatorial verification eq 2}
\widecheck{\beta}(\zeta^m) = \widecheck{\beta}(\beta(x^{-1})) = s_{\beta}(x) \cdot x^{-1}. 
\end{equation}


To check the GKM conditions for $f$, we need to see that if \eqref{equation: combinatorial verification eq 0}, \eqref{equation: combinatorial verification eq 1} hold, then $\loc_{\lambda}(f) - \loc_{\lambda'}(f)$ vanishes at $(x, \zeta)$. We compute
\begin{align*}
\gamma_{\lambda}(f(y)) - \gamma_{\lambda'}(f(y)) 
&= f(\lambda(\zeta) \cdot x) - f(s_{\beta}(\lambda(\zeta)) \cdot \widecheck{\beta}(\zeta^m) \cdot x) \\ &= f(\lambda(\zeta) \cdot x) - f(s_{\beta}(\lambda(\zeta)) \cdot s_{\beta}(x) \cdot x^{-1} \cdot x) \\ 
&= f(\lambda(\zeta) \cdot x) - f(s_{\beta}(\lambda(\zeta) \cdot x)) \\
&= f(\lambda(\zeta) \cdot x) - f(\lambda(\zeta) \cdot x) \\
&= 0
\end{align*}
Here, the first line is the definition of $\gamma$ and \eqref{equation: combinatorial verification eq 0}, the second follows by \eqref{equation: combinatorial verification eq 2}, the third is clear, and the fourth holds since $f$ is $W$-invariant. 
\end{proof}



\subsubsection{Factorization through the deformation.}
We claim that the localization map $\loc$ factors through our deformation.
\begin{lemma}
There is a factorization
\begin{equation*}
\begin{tikzcd}
    \O(T \times \Gm^{\rot} \times T \sslash W) \arrow[r, "\loc"] \arrow[d, hookrightarrow]& \GKM^{0}_{T \times \Gm^{\rot}}(\Gr_G; \k) \arrow[r, hookrightarrow] &  \prod_{\lambda \in X_{\bullet}} \O(T \times \Gm^{\rot})\\
    \O(\eN_{\mu_{\infty}}(\eZ, T \times T \sslash W)) \arrow[ru, swap, dashed, "\bloc"]
\end{tikzcd}
\end{equation*}
\end{lemma}
\begin{proof}
We first check that the whole horizontal composition factors through $\O(\eN_{\mu_{\infty}}(\eZ, T \times T \sslash W))$.
To this end, we need to check that for each root of unity $\zeta$ of $\ord(\zeta) = \ell$ and for every $f \in \O(T \times \Gm^{\rot} \times T \sslash W)$ that vanishes on $Z_{\ell}$, the image $\loc(f)$ inside $\prod_{\lambda \in X_{\bullet}} \O(T \times \Gm^{\rot})$ is divisible by $(q -\zeta)$.
%
This can be checked componentwise -- we need to see that each $\loc_{\lambda}(f) \in \O(T \times \Gm^{\rot})$ is divisible by $(q-\zeta)$, or equivalently that $\loc_{\lambda}(f)$ vanishes after substituing $q=\zeta$. 


To check this, we may identify $f$ with its image under $\O(T \times T\sslash W) \hookrightarrow \O(T \times T)$. For any $(x, \zeta)$, the image of $f$ under the localization map is computed as
\begin{equation*}
\loc_{\lambda}(f) = f(x, \lambda(\zeta) \cdot x).  
\end{equation*}
Specializing to our root of unity $\zeta$, this is the value of $f: T \times T \to \A^1$ on the $T$-valued point $(x, \lambda(\zeta) \cdot x): T \to T \times \Gm^{\rot} \times T$. Since $\zeta$ is an $\ell$-th root of unity, the point $\lambda(\zeta) \in T$ is $\ell$-torsion: we have $\lambda(\zeta)^{\ell} = 1$.
Hence
\begin{equation*}
    x^{\ell} = (\lambda(\zeta)\cdot x)^{\ell}.
\end{equation*}
In other words, $(x, \lambda(\zeta) \cdot x)$ is a $T$-valued point of the closed subscheme $Z'_{\ell} \hookrightarrow T \times T$. Since $f \in \eI_{Z'_{\ell}}$ by assumption, we deduce the desired vanishing $f(x, \lambda(\zeta) \cdot x) = 0$.

To conclude, we need to see that the image of $\O(\eN_{\mu_{\infty}}(\eZ, T \times T \sslash W))$ inside $\prod_{\lambda \in X_{\bullet}} \O(T \times \Gm^{\rot})$ satisfies the GKM conditions. Given any $f \in \eI_{Z_{\ell}'}$ and $\lambda, \lambda' \in X_{\bullet}$ with $ \lambda' = s_{\beta, m}\cdot \lambda$, we thus need to check that the element
\begin{equation*}
\loc_{\lambda}(\tfrac{f}{q - \zeta}) - \loc_{\lambda'}(\tfrac{f}{q - \zeta})
= \tfrac{\loc_{\lambda}(f) - \loc_{\lambda'}(f)}{q - \zeta} \in \O(T \times \Gm^{\rot})
\end{equation*}
is divisible by $(1 - \beta q^m)$.
We already know this is the case for any element in $\O(T \times T \sslash W)$ by Lemma \ref{lemma: GKM localization map combinatorial}. Since $(q - \zeta) \lhd \O(T \times \Gm^{\rot})$ is a prime ideal which does not contain $(1-\beta q^m)$, we deduce that the above ratio is also divisible by $(1 - \beta q^m)$ inside $\O(T \times \Gm^{\rot})$ as desired.
\end{proof}

\subsubsection{Injectivity.}\label{subsubsection: injectivity}
The injectivity of the above constructed map is almost obvious.
\begin{lemma}\label{lemma: injectivity}
The following ring homomorphisms introduced above are injective
\begin{equation*}
    \O(T \times \Gm^{\rot} \times T \sslash W) \to \O(\eN_{\mu_{\infty}}(\eZ, T \times T \sslash W))  \to \GKM^{0}_{T \times \Gm^{\rot}}(\Gr_G; \k) \to \prod_{\lambda \in X_{\bullet}} \O(T \times \Gm^{\rot}).
\end{equation*}
\end{lemma}
\begin{proof}
All four rings in question are clearly reduced (the leftmost one are functions on a reduced scheme; the second one is given by adding certain elements from its fraction field; the third one is a subring of the rightmost infinite product of polynomial rings). Since they are moreover algebras over $\O(T \times \Gm^{\rot})$, they all embed into the localizations at the generic point $\eta \in T \times \Gm^{\rot}$ and it is enough to check the injectivity there. But over $\eta$, the first map becomes an isomorphism (by construction); the second map is injective (it is given by restricting functions on a localization of the affine space $\A^n_{\eta} = \Spec(k(\eta)(T\sslash W))$ over the fraction field $k(\eta)$ to an infinite discrete subset $X_{\bullet} \subseteq \A^n_{\eta}$ that spans a full rank sublattice and hence is Zariski dense); the third map becomes an isomorphism (by construction).
\end{proof}



\subsubsection{Schubert completion.}
Consider the closed immersions \eqref{equation: Gamma and Upsilon}. Taking global functions, we get the maps $\O(\eN_{\mu_{\infty}}(\eZ, T \times T \sslash W)) \to \O(\overline{\Upsilon}) \to \O(\overline{\Upsilon}_{\wt(\mu)})$. The composites $\O(\eN_{\mu_{\infty}}(\eZ, T \times T \sslash W)) \twoheadrightarrow \O(\overline{\Upsilon}_{\wt(\mu)})$ are surjective and compatible as $\mu$ varies. Their kernels induce an exhaustive filtration by ideals $(\eI_{\leq \mu})_{\mu \in X_{\bullet}}$ on $\O(\eN_{\mu_{\infty}}(\eZ, T \times T \sslash W))$. For convenience, we also denote the quotients by $\O(\eN_{\mu_{\infty}}(\eZ, T \times T \sslash W))_{\leq \mu} := \O(\eN_{\mu_{\infty}}(\eZ, T \times T \sslash W)) / \eI_{\leq \mu} $. We obtain the diagram
\begin{equation}\label{equation: schubert completion square}
\begin{tikzcd}
\O(\eN_{\mu_{\infty}}(\eZ, T \times T \sslash W)) \arrow[r, hook] \arrow[d, two heads] & \O(\overline{\Upsilon})  \arrow[d, two heads] \\
\O(\eN_{\mu_{\infty}}(\eZ, T \times T \sslash W))_{\leq \mu} \arrow[r, equal] & \O(\overline{\Upsilon}_{\wt(\mu)}).
\end{tikzcd}
\end{equation}
Here, the upper horizontal map is injective: restricting to irreducible components of $\overline{\Upsilon}$, the ring $\O(\overline{\Upsilon})$ embeds into $\prod_{\lambda \in X_{\bullet}} \O(T \times \Gm^{\rot})$; the induced map from $\O(\eN_{\mu_{\infty}}(\eZ, T \times T \sslash W))$ is precisely the injection from Lemma \ref{lemma: injectivity}. The filtration $(\eI_{\leq \mu})_{\mu \in X_{\bullet}}$ is thus separated.


\begin{definition}\label{definition: Schubert completion}
    The \textit{Schubert completion} of $\O(\eN_{\mu_{\infty}}(\eZ, T \times T \sslash W))$ is the completion with respect to the above separated exhaustive filtration. It is given by
    \begin{equation*}
        \widehat{\O}(\eN_{\mu_{\infty}}(\eZ, T \times T \sslash W)) := \lim_{\mu} {\O}(\eN_{\mu_{\infty}}(\eZ, T \times T \sslash W))_{\leq \mu}.
    \end{equation*}
    Since the filtration is separated, we have a natural injection
    \begin{equation*}
     \O(\eN_{\mu_{\infty}}(\eZ, T \times T \sslash W)) \hookrightarrow  \widehat{\O}(\eN_{\mu_{\infty}}(\eZ, T \times T \sslash W)).
    \end{equation*}
\end{definition}


With the above definitions, the following statement is clear.
\begin{lemma}\label{lemma: completed deformation as functions on Upsilon}
We have a ring isomorphism
    \begin{equation*}
        \widehat{\O}(\eN_{\mu_{\infty}}(\eZ, T \times T \sslash W)) = \O(\overline{\Upsilon}).
    \end{equation*}
\end{lemma}
\begin{proof}
Apply $\lim_{\mu}(-)$ to the bottom row in \eqref{equation: schubert completion square} varying in $\mu$. By definition of the Schubert completion, $\widehat{\O}(\eN_{\mu_{\infty}}(\eZ, T \times T \sslash W)) = \lim_{\mu} \O(\eN_{\mu_{\infty}}(\eZ, T \times T \sslash W))_{\leq \mu}$. Computing functions from an ind-presentation of $\overline{\Upsilon}$, also $\O(\overline{\Upsilon}) = \lim_{\mu} \O(\overline{\Upsilon}_{\wt(\mu)})$.
\end{proof}


We also easily deduce the following.
\begin{lemma}\label{lemma: completed comparison}
The above induces an injective ring homomorphism
\begin{equation}\label{equation: completed comparison}
    \widehat{\O}(\eN_{\mu_{\infty}}(\eZ, T \times T \sslash W)) \hookrightarrow \GKM^{0}_{T \times \Gm^{\rot}}(\Gr_G; \k).
\end{equation}
\end{lemma}
\begin{proof}
To construct the map, observe that $\GKM^{0}_{T \times \Gm^{\rot}}(\Gr_G; \k)$ is complete with respect to the filtration given by kernels of $\GKM^{0}_{T \times \Gm^{\rot}}(\Gr_G; \k) \to \GKM^{0}_{T \times \Gm^{\rot}}(\Gr_{\leq \mu}; \k)$. Under the injective homomorphism ${\O}(\eN_{\mu_{\infty}}(\eZ, T \times T \sslash W)) \hookrightarrow \GKM^{0}_{T \times \Gm^{\rot}}(\Gr_G; \k)$ from Lemma \ref{lemma: injectivity}, this matches the Schubert filtration on the source. Consequently, it induces a map \eqref{equation: completed comparison}. This map is injective as both sided embed compatibly into $\prod_{\lambda \in X_{\bullet}} \O(T \times \Gm^{\rot})$.
\end{proof}



\subsection{The topological meaning}\label{section: map to equivariant K-theory: topology}
We now explain the topological meaning of the map from \S \ref{section: map to equivariant K-theory: combinatorics}. This ties to the previous section by localization to fixed points. Ignoring this relationship, the topological arguments are completely independent.
In particular, we interpret the map $\loc$ and fundamental Bezrukavnikov classes in purely topological terms. This interpretation makes sense over any coefficient ring $\k$, for example over $\Z$.

\subsubsection{The map from functions on tori.}\label{section: the map from functions on tori}
We work with a reductive group $G$.
To allows for topological arguments, the geometry of $\Gr_G$ takes place over $k=\C$. However, the coefficient ring $\k$ may be taken arbitrary.

There is a natural map
\begin{equation*}
 K^{\top}_{T \times \Gm^{\rot} \times G}(\pt) \xrightarrow{} K^{\top}_{T \times \Gm^{\rot}}(\Gr_G)  
\end{equation*}
defined as the equivariant structure map under the vertical identifications
\begin{equation*}
    \begin{tikzcd}[row sep=5]
      K^{\top}_{T \times \Gm^{\rot} \times G}(\pt) \arrow[r] & K^{\top}_{T \times \Gm^{\rot}}(\Gr_G) \\
      & K^{\top}((T \times \Gm^{\rot}) \backslash \lo G / \lop G) \arrow[u, equal] \\
      K^{\top}((T \times \Gm^{\rot}) \backslash \pt / G) \arrow[uu, equal] \arrow[r] & K^{\top}((T \times \Gm^{\rot}) \backslash \lo G / G) \arrow[u, equal].
    \end{tikzcd}
\end{equation*}

\begin{remark}
More generally, given two parahoric subgroups $\eP$ and $\eP'$ of $\lo G$ with associated Levi quotients $L$ and $L'$, we get the map
\begin{equation*}
 K^{\top}_{L \times \Gm^{\rot} \times L'}(\pt) \to K^{\top}_{\eP \rtimes \Gm^{\rot}}(\lo G/ \eP').  
\end{equation*}    
\end{remark}

\subsubsection{Localization map.}
We can now identify
\begin{equation}\label{equation: localization map}
   \begin{tikzcd}[row sep=5]
        K^{\top, 0}_{T \times \Gm^{\rot} \times G}(\pt; \k) \arrow[r] & K^{\top, 0}_{T \times \Gm^{\rot}}(\Gr_G; \k) \arrow[r, hookrightarrow] & K^{\top, 0}_{T \times \Gm^{\rot}}(\Gr^{T \times \Gm^{\rot}}; \k) \\
        \O(T \times \Gm^{\rot} \times T \sslash W) \arrow[r] \arrow[u, equal] & \GKM^{0}_{T \times \Gm^{\rot}}(\Gr_G; \k) \arrow[u, equal] \arrow[r, hookrightarrow] & \prod_{\lambda \in X_{\bullet}} \O(T \times \Gm^{\rot}) \arrow[u, equal].
    \end{tikzcd}
\end{equation} 
In particular, this gives a canonical map
\begin{equation}\label{equation: map from two copies of tori}
 \O(T \times \Gm^{\rot} \times T \sslash W) \to K^{\top, 0}_{T \times \Gm^{\rot} \times G}(\Gr_G; \k)
\end{equation}
Under the identification from \eqref{equation: localization map}, the left-right composition is the localization map $\gamma$ from \eqref{equation: combinatorial localization map}. 


\subsubsection{Tautological bundles on $\Gr$.}
Let us first take $G=GL_n$. Let $\Lambda_0 = k\llb t \rrb^{\oplus n} \subseteq k\llp t \rrp^{\oplus n}$ be the standard lattice. Fix some $\mu \in X_{\bullet}$ and let $\Gr_{\leq \mu}$ be the corresponding affine Schubert variety. It parametrizes certain sublattices $\Lambda \subseteq k\llp t \rrp^{\oplus n}$. Depending on $\mu$, we choose a sufficiently big positive integer $m \geq 0$ such that for each $\Lambda \in \Gr_{\leq \mu}$ it holds that $t^m \Lambda_0 \subseteq \Lambda \subseteq t^{-m} \Lambda_0$.

After fixing this normalization, we have the tautological vector bundle $\eE_{\leq \mu}$ on $\Gr_{\leq \mu}$, associating to an $R$-valued points $\Lambda$ of $\Gr_{\leq \mu}$ the algebraic vector bundle $t^{-m}\Lambda_0 / \Lambda$ on $R$:
\begin{equation*}
    \eE_{\leq \mu}: \Lambda \mapsto (\coker \Lambda \to t^{-m}\Lambda_0).
\end{equation*}
This is indeed a vector bundle, see \cite[proof of Theorem 1.1.3]{Zhu15}.
Similarly, we have the trivial rank $mn$ vector bundle $\eG_{\leq \mu}$ given by 
$\coker (\Lambda_0 \to t^{-m}\Lambda_0)$.
Both of these are naturally equivariant with respect to the groups $T \times \Gm^{\rot} \leq \lop GL_n \rtimes \Gm^{\rot}$.


\subsubsection{Bezrukavnikov classes for $GL_n$ from tautological bundles.}\label{section: bezrukavnikov classes for GL_n from tautological bundles}
Let $G=GL_n$. Consider the equivariant $K$-theory class given by the difference of the above vector bundles
\begin{equation}\label{equation: normalized tautological class 1}
[\eE_{\leq \mu}] - [\eG_{\leq \mu}] = [t^{-m}\Lambda_0 / \Lambda] - [t^{-m}\Lambda_0 / \Lambda_0] \in K^{\top, 0}_{T \times \Gm^{\rot}}(\Gr_{\leq \mu}; \k)    
\end{equation}
This can be rewritten in several ways. Firstly, it is given by the perfect complex
\begin{equation}\label{equation: normalized tautological class 2}
[\Lambda \to t^{-m}\Lambda_0] + [\Lambda_0 \to t^{-m}\Lambda_0][1]
\end{equation}
on $\Gr_{\leq \mu}$; here the $t^{-m}\Lambda_0$ sits in degree zero and $[1]$ is a homological shift. 
%
Alternatively, it can be also rewritten as a difference of classes of equivariant coherent sheaves which lie in $\Perf(\Gr_{\leq \mu})$, namely
\begin{equation}\label{equation: normalized tautological class 3}
    [\Lambda_0 / \Lambda_0 \cap \Lambda] - [\Lambda / \Lambda_0 \cap \Lambda].
\end{equation}
Informally, it may be thought of as the difference $[\Lambda_0] - [\Lambda]$, but this does not quite make sense since neither of these projective modules is of finite rank.

In any case, the class \eqref{equation: normalized tautological class 1} is clearly independent of the choice of $m$, and compatible with varying $\mu$. (This is obvious from \eqref{equation: normalized tautological class 3}; it alternatively follows by $[t^{-m}\Lambda_0 / \Lambda] - [t^{-m}\Lambda_0 / \Lambda_0] = [t^{-m}\Lambda_0 / \Lambda] + [t^{-m'}\Lambda_0 / t^{-m}\Lambda_0 ] - [t^{-m'}\Lambda_0 / t^{-m}\Lambda_0 ] - [t^{-m}\Lambda_0 / \Lambda_0] = [t^{-m'}\Lambda_0 / \Lambda] - [t^{-m'}\Lambda_0 / \Lambda_0]$ for any $m' \geq m$.)  It thus gives a well defined equivariant $K$-theory class on the whole of $\Gr_G$. We then have the following.

\begin{comment}
\begin{remark}
Note that due to renormalization issues, $\eE_{\leq \mu}$ do not assemble into a single vector bundle on $\Gr$. Indeed, as $\mu$ varies, so does $\rank \eE_{\leq \mu}$. However, the top exterior power $\wedge^{\rank \eE_{\leq \mu}} \eE_{\leq \mu}$ is a well-defined line bundle on $\Gr_G$, called the \textit{determinant line bundle} $\eL_{\det}$.
\end{remark}
\end{comment}


\begin{lemma}
Inside $K^{\top, 0}_{T \times \Gm^{\rot}}(\Gr_G; \k)$, we have $b_{1,1} = [t^{-m}\Lambda_0 / \Lambda] - [t^{-m}\Lambda_0 / \Lambda_0]$.
\end{lemma}
\begin{proof}
Given any $\lambda \in X_{\bullet}$, it is easy to compute the localizations of both sides to the corresponding fixed-point. Denote $\lambda = (t_1^{d_1}, \dots, t_n^{d_n})$ for $d_1, \dots, d_n \in \Z$.
%
%Given any $d \in \Z$, we denote its sign by $\epsilon(d) \in \{-1, 0, 1\}$.
For any $d \in \Z$, note that
\begin{equation}
[d]_q := \frac{1-q^d}{1-q} =   
\begin{cases}\label{equation: cases for geometric series}
1 + q + \dots + q^{d-1} & \text{if $d > 0$,} \\
0 & \text{if $d = 0$,} \\
-q^{-1}( 1 + q^{-1} + \dots + q^{d+1}) & \text{if $d < 0$.}
\end{cases}
\end{equation}
We now compute
\begin{align*}
\loc_{\lambda}([t^{-m}\Lambda_0 / \Lambda] - [t^{-m}\Lambda_0 / \Lambda_0])
&= \loc_{\lambda}([\Lambda_0 / \Lambda_0 \cap \Lambda] - [\Lambda / \Lambda_0 \cap \Lambda]) \\
&= \sum_{\{i \mid d_i > 0\}} (1+q+ \dots + q^{d_i-1})t_i - \sum_{\{i \mid d_i < 0\}} (q^{-1}+ \dots + q^{d_i+1})t_i \\
&= \sum_{i=1}^n \frac{1-q^{d_1}}{1-q} \cdot t_i \\
&= \frac{\sum_{i=1}^n t_i - \sum_{i=1}^n q^{d_i} \cdot t_i}{1-q} \\
& =\gamma_{\lambda} \left( \frac{e_1(t_1, \dots, t_n) - e_1(s_1, \dots, s_n)}{1-q} \right)\\
&=\gamma_{\lambda}(b_{1,1}) 
\end{align*}
Here, the first line is \eqref{equation: normalized tautological class 3}. The second line is by concrete knowledge of the fibers of $t^{-m}\Lambda_0 / \Lambda$ and $t^{-m}\Lambda_0 / \Lambda$ at the fixed point $\lambda$ as $T \times \Gm^{\rot}$ representations. The third line is by \eqref{equation: cases for geometric series}. The fourth line is clear, the fifth line follows by the description of $\gamma_{\lambda}$ from \eqref{equation: combinatorial localization map}, and the sixth line is by definition of $b_{1,1}$.

Since the above holds for any $\lambda \in X_{\bullet}$, the lemma follows.
\end{proof}

\subsubsection{Bezrukavnikov classes from vector bundles for general $G$.}
Let $G$ be any reductive group.
The affine Grassmannian is functorial in $G$, as is our description of its equivariant $K$-theory. Consequently, we can define the Bezrukavnikov classes from the standard one on $\Gr_{GL_n}$.

Namely, consider any representation $\rho: G \to GL_n$. This induces a map $\Gr_G \to \Gr_{GL_n}$, equivariant with respect to $\lo G \rtimes \Gm^{\rot} \to \lo GL_n \rtimes \Gm^{\rot}$ and compatible with the Schubert stratifications.
We thus have a pullback map $\rho^*: K^{\top, 0}_{T_{GL_n} \times \Gm^{\rot}}(\Gr_{GL_n}; \k) \to K_{T_{G} \times \Gm^{\rot}}^{\top, 0}(\Gr_G; \k)$. Using the double-quotient description \S \ref{section: local quantum hecke stacks} and the construction \S \ref{section: the map from functions on tori}, it follows that $\rho^*$ induces the obvious map $\rho^*: T_{GL_n} \times \Gm^{\rot} \times T_{GL_n} \to T_{G} \times \Gm^{\rot} \times T_{G}$, and hence on the deformation constructions.


In particular, for any $f \in \eI_{Z_{\ell}} \subseteq \O(T_{GL_n} \times T_{GL_n} \sslash W)$, this map sends $\tfrac{f}{1-q^{\ell}} \mapsto \tfrac{\rho^*(f)}{1-q^{\ell}}$. 
Specializing further to the case when $G=G^{\sc}$ is simply connected, $\rho = \omega_j$ is the $j$-th fundamental representation and $f=e_1(t_1, \dots, t_n) - e_1(s_1, \dots, s_n)$, we deduce it sends $b^{GL_n}_{1,1} \mapsto b^G_{j, 1}$. In other words, $b^G_{j, 1}$ is the equivariant $K$-theory class $\omega_j^*([\Lambda_0 / \Lambda_0 \cap \Lambda] - [\Lambda / \Lambda_0 \cap \Lambda] )$ on $\Gr_{G}$. We have proved the following.

\begin{lemma}\label{lemma: fundamental bezrukavnikov classes as pullbacks} We have
 \begin{equation*}
    b^G_{j,1}= \omega_j^{*} (b^{GL_n}_{1,1}). 
 \end{equation*}   
\end{lemma}

\subsubsection{Adams operations.}
The fundamental Bezrukavnikov classes 
$b_{j, \ell} = \tfrac{e_j(t^{\ell}_1, \dots, t_n^{\ell}) - e_j(s^{\ell}_1, \dots, s_n^{\ell})}{1-q^{\ell}}$
are related to each other through Adams operations in equivariant $K$-theory. 
\begin{lemma}\label{lemma: adams operations on bezrukavnikov classes}
Let $1 \leq j \leq n$. For each $\ell, \ell' \in \N$, we have
\begin{equation}\label{equation: adams operations on bezrukavnikov classes}
    \psi^{\ell}(b_{j, \ell'}) = b_{j, \ell\ell'}.
\end{equation}    
\end{lemma}
\begin{proof}
The Adams operation $\psi^{\ell}$ is a functorial ring endomorphism of $K^{\top, 0}_{T \times \Gm^{\rot}}(-)$, which exponentiate classes of line bundles to the $\ell$-th power.
By injectivity of localization to fixed-points (valid in our setup), it is thus enough to check the equality \eqref{equation: adams operations on bezrukavnikov classes} after restriction to each fixed point, in other words after applying $\gamma_{\lambda}$. 

After this restriction, each of the generators $t_i, s_i, q, 1$ is a $K$-theory class of an equivariant line bundle (a character of $T \times \Gm^{\rot}$). Therefore, $\psi^{\ell}$ acts by $(-)^{\ell}$ on each of them. Since $\psi^{\ell}$ is a ring homomorphism, the desired equality follows.  
\end{proof}






\subsubsection{Topological argument.}
With the above topological meaning at hand, we have the following clear construction, which we again state for $G$ simply connected or $GL_n$.
\begin{lemma}
    We have ring homomorphisms
    \begin{equation*}
        \O(T \times \Gm^{\rot} \times T \sslash W) \to \O(\eN_{\mu_{\infty}}(\eZ, T \times T \sslash W)) \to \widehat{\O}(\eN_{\mu_{\infty}}(\eZ, T \times T \sslash W)) \to K^{\top, 0}_{T \times \Gm^{\rot}}(\Gr_G; \k).
    \end{equation*}
\end{lemma}
\begin{proof}
The map from the leftmost term is given by the equivariant structure maps \S \ref{section: the map from functions on tori}. To see that it factors through the second term, we need to prove that the fundamental Bezrukavnikov classes $b_{j, \ell}$ give well-defined classes in $K$-theory. However, we have already interpreted them as Adams operations on the fundamental vector bundles in Lemmata \ref{lemma: fundamental bezrukavnikov classes as pullbacks}, \ref{lemma: adams operations on bezrukavnikov classes}.

Finally, $K^{\top, 0}_{T \times \Gm^{\rot}}(\Gr_G; \k)$ is complete with respect to the filtration given by kernels of restriction maps to affine Schubert varieties by Proposition \ref{proposition: completion of K-theory}. Equipping $\O(\eN_{\mu_{\infty}}(\eZ, T \times T \sslash W))$ with the Schubert filtration of Definition \ref{definition: Schubert completion} makes the comparison map filtered. Hence it factors through the Schubert completion $\widehat{\O}(\eN_{\mu_{\infty}}(\eZ, T \times T \sslash W))$ as desired.
\end{proof}





\subsection{Bezrukavnikov--Finkelberg description of equivariant \texorpdfstring{$K$}{K}-theory}\label{section: bezrukavnikov--finkelberg description of equivariant K-theory}
We now conclude the proof of our main result for simply connected $G$ with $\C$-coefficients. To this end, we will assume $k = \C = \k$. We further record the topological Adams generation. We find this case particularly illustrative; technical generalizations are then discusses in \S \ref{section: variants and generalizations}. 

\subsubsection{The isomorphism.}
We now prove the surjectivity by a topological argument.
\begin{theorem}\label{theorem: equivariant K-theory as deformation to normal cone -- G sssc}
Assume $G$ is simply-connected group over $\C$. Then the above constructed map induces a ring isomorphism
\begin{equation*}
   \widehat{\O}(\eN_{\mu_{\infty}}(\eZ, T \times T \sslash W)) \xrightarrow{\cong} K^{\top, 0}_{T \times \Gm^{\rot}}(\Gr_G; \C).
\end{equation*}
\end{theorem}
\begin{proof}
Regard both sides as quasi-coherent (so in particular fpqc) sheaves on $T \times \Gm^{\rot}$. To see that the discussed map is an isomorphism, it is thus enough to check it is an isomorphism after formal completion at each closed point $(x, \zeta) \in T \times \Gm^{\rot}$. (Alternatively, it is enough to check the isomorphisms at actual fibers by Remark \ref{remark: replacing formally completed fibers by actual fibers}.) We have thus reduced to checking that for each $(x, \zeta)$, the bottom horizontal map in the following diagram is an isomorphism.
\begin{equation*}
    \begin{tikzcd}
         \widehat{\O}(\eN_{q=\zeta}(\eZ, T \times T \sslash W)) \arrow[r] \arrow[d] & K^{\top, 0}_{T \times \Gm^{\rot}}(\Gr_G; \C) \arrow[d] \\
         \widehat{\O}(\eN_{q=\zeta}(\eZ, T \times T \sslash W))^{\wedge}_{(x, \zeta)} \arrow[r] & K^{\top, 0}_{T \times \Gm^{\rot}}(\Gr_G; \C)^{\wedge}_{(x, \zeta)}
    \end{tikzcd}
\end{equation*}
To this end, fix a closed point $(x, \zeta) \in T \times \Gm^{\rot}$. To check the isomorphism between the completions, we distinguish cases depending on the genericity of $\zeta$.

\bigskip
Assume first that $\zeta$ is a primitive root of unity of order $\ell$. Then consider the following diagram.
\begin{equation*}
    \begin{tikzcd}[row sep = 5, column sep =15]
      \widehat{\O}({\eN}_{\mu_{\infty}}(\eZ, T \times T \sslash W))^{\wedge}_{(x, \zeta)} \arrow[d, equal] \arrow[r] & K^{\top, 0}_{T \times \Gm^{\rot}}(\Gr_{G}; \C)^{\wedge}_{(x, \zeta)} \arrow[d, equal] \\
      %
      \widehat{\O}({N}_{q=\zeta}(Z_{\ell}, T \times T \sslash W))^{\wedge}_{(x, \zeta)} \arrow[d, equal]  & \widehat{H}^{\bullet}_{T \times \Gm^{\rot}}(\Fix_{(x, \zeta)}(\Gr_G); \C)^{\wedge}_{(0, 0)} \arrow[d, equal] \\
      %
      \prod_{\vartheta \in |\overline{\Gamma}_{(x, \zeta)}|} \widehat{\O}(N_{q=\zeta}({Z_{\ell}^{\vartheta}}, T \times T\sslash W))^{\wedge}_{(x, \zeta)} \arrow[d, equal] & \widehat{H}^{\bullet}_{T \times \Gm^{\rot}}(\coprod_{\vartheta \in |\overline{\Gamma}_{(x, \zeta)}|} \Fl^{\vartheta}_{\ell, L^{\sc}_{[x, \zeta]}}; \C)^{\wedge}_{(0, 0)} \arrow[d, equal] \\
      %
      \prod_{\vartheta \in |\overline{\Gamma}_{(x, \zeta)}|} \widehat{\O}(N_{h=0}(\t \times_{\t \sslash W_{[x, \zeta]}} \t \sslash W_{\vartheta},  \t \times \t \sslash W_{\vartheta} ))^{\wedge}_{(0,0)} \arrow[r, equal] & \prod_{\vartheta \in |\overline{\Gamma}_{(x, \zeta)}|} \widehat{H}^{\bullet}_{T \times \Gm^{\rot}}(\Fl^{\vartheta}_{\ell, L^{\sc}_{[x, \zeta]}}; \C)^{\wedge}_{(0, 0)}
    \end{tikzcd}
\end{equation*}
The identifications in the left-hand column work as follows. First, since we are completing at $(x, \zeta)$, only the deformation to the normal cone of $Z_{\ell}$ centered at $\zeta$ matters. The next two identifications are covered by Lemma \ref{lemma: completions of the deformation}.

The identifications in the right-hand column go a follows. First use the localization from Corollary \ref{corollary: K-theory in the flat situation}, Theorem \ref{theorem: fibers of topological K-theory, HP and fixed-points} and Remark \ref{remark: the completed variant of equivariant localization} to pass from the completion of equivariant $K$-theory over $(x, \zeta)$ to the completion of equivariant cohomology of the fixed-points of $(x, \zeta)$ at the origin (alternatively see Remark \ref{remark: replacing formally completed fibers by actual fibers}). Then use the computation of the reduced fixed-points on the affine Grassmannian from Theorem \ref{lemma: fixed points of (c, zeta) on Gr_G -- zeta root of unity}. Finally, we may commute the disjoint union through cohomology.

The bottom identification is by taking the Schubert completion of the isomorphism from Proposition \ref{proposition: cohomology of affine flag varieties as deformation to the normal cone}. (Note that the completion of equivariant cohomology is given by the Schubert completion.) This finishes the discussion of the case when $\zeta$ is a root of unity.

\bigskip
Now, assume on the other hand that $\zeta \in \Gm^{\rot} - \mu_{\infty}$. We then make the analogous identifications.
\begin{equation*}
    \begin{tikzcd}[row sep = 5, column sep =15]
      \widehat{\O}({\eN}_{\mu_{\infty}}(\eZ, T \times T \sslash W))^{\wedge}_{(x, \zeta)} \arrow[d, equal] \arrow[r] & K^{\top, 0}_{T \times \Gm^{\rot}}(\Gr_{G}; \C)^{\wedge}_{(x, \zeta)} \arrow[d, equal] \\
      %
      \widehat{\O}(T \times \Gm^{\rot} \times T \sslash W)^{\wedge}_{(x, \zeta)} \arrow[d, equal]  & \widehat{H}^{\bullet}_{T \times \Gm^{\rot}}(\Fix_{(x, \zeta)}(\Gr_G); \C)^{\wedge}_{(0, 0)} \arrow[d, equal] \\
      %
      \prod_{\vartheta \in |\overline{\Gamma}_{(x, \zeta)}|} {\O}(\Gm^{\rot} \times T \times_{T \sslash W_{[x, \zeta]}} T \sslash W_{\vartheta} )^{\wedge}_{(x, \zeta)} \arrow[d, equal] & {H}^{\bullet}_{T \times \Gm^{\rot}}(\coprod_{\vartheta \in |\overline{\Gamma}_{(x, \zeta)}|} L_{[x, \zeta]}/P_{\vartheta}; \C)^{\wedge}_{(0, 0)} \arrow[d, equal] \\
      %
      \prod_{\vartheta \in |\overline{\Gamma}_{(x, \zeta)}|} {\O}(\Ga \times\t \times_{\t \sslash W_{[x, \zeta]}} \t \sslash W_{\vartheta})^{\wedge}_{(0,0)} \arrow[r, equal] & \prod_{\vartheta \in |\overline{\Gamma}_{(x, \zeta)}|} {H}^{\bullet}_{T \times \Gm^{\rot}}(L_{[x, \zeta]}/P_{\vartheta}; \C)^{\wedge}_{(0, 0)}
    \end{tikzcd}
\end{equation*}

The identifications in the left-hand column are as follows. Firstly, since $\zeta \notin \mu_{\infty}$, after the completion at $(x, \zeta)$ the deformations to the normal cones do not play a role. For the second identification, note that the Schubert completion corresponds to a limit over restrictions to finite subsets of the component set $|\Gamma_{(x, \zeta)}|$. Each of the connected components $\Gamma^{\vartheta}_{(x, \zeta)}$, $\vartheta \in |\Gamma_{(x, \zeta)}|$ is given by $W_{[x, \zeta]}$ copies of $T$ intersection in the natural way; the stabilizer of this subscheme with respect to the $W$-action on the right copy of $T$ is precisely $W_{\vartheta}$. The third identification is the usual comparison of the completion of a semisimple, simply-connected group with its Lie algebra.

For the right hand column, the first identification is the localization of equivariant $K$-theory as in Corollary \ref{corollary: K-theory in the flat situation}, Theorem \ref{theorem: fibers of topological K-theory, HP and fixed-points} and Remark \ref{remark: the completed variant of equivariant localization}. The second isomorphism is the computation of reduced fixed points from Theorem \ref{lemma: fixed points of (c, zeta) on Gr_G -- zeta generic} (no completion is necessary, since each of the terms has bounded cohomology). The third identification is clear.

The bottom horizontal identification is the well-known description of $T$-equivariant cohomology of partial flag varieties of reductive groups from Lemma \ref{lemma: equivariant cohomology of finite flag varieties} and Remark \ref{remark: finite flag varieties with trivial loop rotation}; note that $\Gm^{\rot}$ acts trivially on each of them.
\end{proof}


\begin{remark}\label{remark: replacing formally completed fibers by actual fibers} 
In fact, in the first reduction step, we may reduce further.
Firstly, the map in question is injective by \S \ref{subsubsection: injectivity}, so we only need to check surjectivity. 

Secondly, since we are speaking about a map of quasi-coherent sheaves on an affine scheme, it is enough to check that for each closed point $(x, \zeta)$, the following map between the fibers is surjective:
\begin{equation*}
    \begin{tikzcd}
        \widehat{\O}(\eN_{q=\zeta}(\eZ, T \times T \sslash W))_{(x, \zeta)} \arrow[r] & K^{\top, 0}_{T \times \Gm^{\rot}}(\Gr_G; \C)_{(x, \zeta)}.
    \end{tikzcd}
\end{equation*}    
Consequently, all occurrences of the formal completion $(-)^{\wedge}_{(x, \zeta)}$ may be replaced by the actual fiber $(-)_{(x, \zeta)}$ throughout the proof. This variant of the proof thus needs a slightly weaker technical input from \S \ref{section: K-theory and HP and cohomology of fixed points}, which can be covered by more standard results.
\end{remark}

\begin{remark}
The horizontal identifications in the proof may be also compared through the GKM presentations, as outlined in Remarks \ref{remark: GKM for cohomology of finite flag varieties}, \ref{remark: informal GKM localization}, \ref{remark: Upsilon and GKM}. This is another sanity check for the compatibility of all the identifications that we made.

Since the GKM presentation works over more general coefficient rings $\k$, this variant of the argument is applicable. In particular, the variant of Theorem \ref{theorem: equivariant K-theory as deformation to normal cone -- G sssc} with coefficients in $\k = \Q$ holds without without modifications. On the other hand, for $\k=\Z$, further blowups over all primes $p \in \Z$ are necessary; we will discuss this elsewhere.
\end{remark}



\subsubsection{Finite generation as complete topological $\lambda$-ring.}
As we saw, the equivariant $K$-theory ring of the affine Grassmannian is quite big -- it is a topological completion of an already non-noetherian ring of functions on our affine blowup. Nevertheless, it is finitely generated in the following appropriate sense.
\begin{corollary}\label{corollary: topological Adams generation}
Let $G$ be a simply connected group over $\C$. As complete topological $\lambda$-algebra over the equivariant base, $K^{\top, 0}_{T \times \Gm^{\rot}}(\Gr_G; \C)$ is generated by the finitely many elements
    \begin{equation*}
        (b_{j, 1} \mid j = 1, \dots, n).
    \end{equation*}
\end{corollary}
\begin{proof}
By Theorem \ref{theorem: equivariant K-theory as deformation to normal cone -- G sssc} and Lemma \ref{lemma: fundamental bezrukavnikov classes}, we know that $K^{\top, 0}_{T \times \Gm^{\rot}}(\Gr_G; \C)$ is topologically generated by the fundamental classes $b_{j, \ell}$. Under $\lambda$-operations, these are in turn generated by $b_{j, 1}$ for $j=1, \dots, n$ by Lemma \ref{lemma: adams operations on bezrukavnikov classes}.
\end{proof}

\subsubsection{Some intuition.}
One could try to approach the surjectivity in the proof of Theorem \ref{theorem: equivariant K-theory as deformation to normal cone -- G sssc} in another way as follows. By Lemma \ref{lemma: completed deformation as functions on Upsilon} and Proposition \ref{proposition: GKM description of K-theory of Gr} we want to prove that the following map is an isomorphism
\begin{equation*}
   \O(\overline{\Upsilon}) \to \GKM^{0}_{T \times \Gm^{\rot}}(\Gr_G; \k). 
\end{equation*}
Both sides compatibly embed into
\begin{equation}\label{equation: product of copies of torus}
    \prod_{\lambda \in X_{\bullet}} \O(T \times \Gm^{\rot}).
\end{equation}
and we are left to check that the images agree. For the right hand side, the image is given by the GKM conditions, saying that functions on different copies of $T \times \Gm^{\rot}$ agree over subtori of the form $T_{\beta + m\hbar} = \ker(\beta + m\hbar)$.
For the left-hand side, we are considering functions on an arrangement of copies of $T \times \Gm^{\rot}$ inside $T \times \Gm^{\rot} \times T\sslash W$. Their intersections combinatorially match the subvarieties imposing the equalizer conditions in the GKM picture.

From what we already said, one would be tempted to conclude. However, this is not a complete reasoning, as Example \ref{example: seminormality issues} below shows. To complete the above intuition, one needs to, in one way or another, study precisely how $\overline{\Upsilon}$ sits inside $\Bl^{\circ}_{\eZ}( T \times \Gm^{\rot} \times T \sslash W)$ via its defining embedding -- this is what the above proof accomplishes.

\begin{example}\label{example: seminormality issues}
 Consider the following example. Let $X_3 := \Spec \k[x,y, z] / (xy, yz, zx)$ be three lines in general position inside the three-space $\A^3$ intersecting in the origin. Let $X_2 := \Spec \k[x, y] / xy(x-y)$ be three lines inside the plane $\A^2$ intersection in the origin.  There is a natural squishing map $X_3 \to X_2$. This map is a homeomorphism -- in fact the same arrangement of lines. However, it is not an isomorphism of schemes: the induced injection on the rings of functions $\O(X_2) \to \O(X_3)$ is not surjective. The problem is related to the notion of \textit{semi-normality} ($X_3$ is seminormal but $X_2$ is not), or even more explicitly to that of \textit{embedding dimension} (which is $3$ and $2$, respectively).   
\end{example}



\subsection{Variants and generalizations}\label{section: variants and generalizations}

\subsubsection{The general answer.}
We work with coefficients $\k=\C$. As stated in Remark \ref{remark: variant of the K-theory of affine Grassmannian}, our Theorem \ref{theorem: equivariant K-theory as deformation to normal cone -- G sssc} has several direct generalizations. These may be either proved by a direct modifications of the above argument, or bootstrapped from what we already have. Putting these remarks together, the generalized statement is as follows. Recall that $G^{\sc}$ denotes the simply connected cover of the derived group of $G$ from \S \ref{subsubsection: simply connected covers and adjoint groups}. We use the modified family of blowup centers $\widetilde{\eZ}$ from \S \ref{subsubsection: modifications in the reductive case}, \S \ref{section: action of the fundamental group}.

\begin{theorem}\label{theorem: equivariant K-theory as deformation to normal cone -- general case}
Let $G$ be a reductive group over $\C$. Denote by $T^{\sc}$ the maximal torus of $G^{\sc}$. Pick two parahoric subsgroups $\eP_1, \eP_2 \leq \lo G$. Consider the affine flag variety $\Fl^J_G := \lo G / \eP_2$, together with the action of $\eP_1 \rtimes \Gm^{\rot} \acts \Fl^J_G$.
Then 
\begin{equation*}
\widehat{\O}(\Bl^{\circ}_{\widetilde{\eZ}}(\Gm^{\rot} \times T\sslash W_1 \times T^{\sc} \sslash W_2 \times \pi_1(G)) = K^{\top, 0}_{\eP_1 \rtimes \Gm^{\rot}}(\Fl^J_G, \C).
\end{equation*}
\end{theorem}


\begin{proof}
The setup of theorem has four degrees of freedom -- we can independently choose:
(i) which parahoric will act on the left;
(ii) which partial affine flag variety are we taking;
(iii) which isogeny class of $G$ we take for $\lo G$, or more precisely which connected components of $\Gr_{G^{\ad}}$ we take into account; 
(iv) which isogeny class of $G$ will be used for the $\lop G$-action, or alternatively which torus $T$ with $T^{\der}$ in between $T^{\sc} \to T^{\der} \to T^{\ad}$ will act.


We will treat these independently: (i) and (ii) in Lemma \ref{lemma: equivariant K-theory as deformation to normal cone -- flag varieties and equivariance}; (iii) in Lemma \ref{lemma: equivariant K-theory as deformation to normal cone -- G reductive}; (iv) in Lemma \ref{lemma: equivariant K-theory as deformation to normal cone -- isogeneous torus} and Remark \ref{remark: equivariant K-theory as deformation to normal cone -- isogeneous torus}. It is quite clear that all this bookkeeping can be done simultaneously; we leave the details to the reader. 
\end{proof}

In what follows, we also write $\Bl^{\circ}_{\widetilde{\eZ}}(X)$ for the affine blowup at the relevant quotients and restriction of the family $\widetilde{\eZ}$ which lives in the given variant $X$ of the scheme $\Gm^{\rot} \times T\sslash W_1 \times T^{\sc} \sslash W_2 \times \pi_1(G)$. This should not lead to any confusion.

\subsubsection{Affine flag varieties and positive loop-group equivariance.}
Assume $G = G^{\sc}$ is simply connected and consider its torus $T=T^{\sc}$. Take two parahoric subgroups $\eP_1, \eP_2 \leq \lo G$. Consider the corresponding affine flag variety $\Fl^J = \lo G / \eP_2$, together with the action
$\eP_1 \rtimes \Gm^{\rot} \acts  \Fl^J_G$. The equivariant $K$-theory of this action is the $K$-theory of the double quotient (quantum Hecke stack)
\begin{equation*}
    \eP_1 \rtimes \Gm^{\rot} \backslash \lo G \rtimes \Gm^{\rot} / \eP_2 \rtimes \Gm^{\rot}
\end{equation*}
which has an apparent symmetry. 

\begin{lemma}\label{lemma: equivariant K-theory as deformation to normal cone -- flag varieties and equivariance}
For $G$ simply connected, we have a ring isomorphism 
\begin{equation*}
    \widehat{\O}(\Bl^{\circ}_{\widetilde{\eZ}}(T\sslash W_1 \times \Gm^{\rot} \times T\sslash W_2))
    =K^{\top, 0}_{\eP_2 \rtimes \Gm^{\rot}}(\Fl^J_G; \C).
\end{equation*}   
\end{lemma}
\begin{proof}
By homotopy invariance and equivariant formality, we have
\begin{equation*}
K^{\top, 0}_{\eP_1 \rtimes \Gm^{\rot}}(\Fl^J_G; \C) = K^{\top, 0}_{L_1 \times \Gm^{\rot}}(\Fl^J_G; \C) = (K^{\top, 0}_{T \times \Gm^{\rot}}(\Fl^J_G; \C))^{W_1}.    
\end{equation*}
Here, the $W_1$-action is on the first copy of the torus. 

For the other part, note that $\Fl^J_G \to \Gr_G$ is a fpqc locally trivial fibration with fiber $G/P_{2}$. The $K$-theory of $\Fl^J_G$ is thus given by the following base change:
\begin{equation*}
    K^{\top, 0}_{T\times \Gm^{\rot}}(\Fl^J_G; \C) = K^{\top, 0}_{T\times \Gm^{\rot}}(\Gr_G; \C) \otimes_{\O(T \sslash W)} \O(T \sslash W_2).
\end{equation*}
Since blowups are compatible with base change, we have reduced the desired statement to Theorem \ref{theorem: equivariant K-theory as deformation to normal cone -- G sssc}.(Alternatively, one can directly repeat our computation of $T \times \Gm^{\rot}$-equivariant $K$-theory of $\Gr$ in the more general case of $\Fl^J_G$.)
\end{proof}


\subsubsection{Affine Grassmannian of reductive $G$.}
Assume now $G$ to be merely reductive, but let us work equivariantly with respect to the torus $T^{\sc}$ of $G^{\sc}$, see \S \ref{subsubsection: simply connected covers and adjoint groups}.

\begin{lemma}\label{lemma: equivariant K-theory as deformation to normal cone -- G reductive}
Assume $G$ is connected reductive. Then the above constructed map induces a ring isomorphism
\begin{equation*}
   \widehat{\O}(\Bl^{\circ}_{\widetilde{\eZ}}(T^{\sc} \times \Gm^{\rot} \times T^{\sc}\sslash W \times \pi_1(G))) 
   = K^{\top, 0}_{T^{\sc} \times \Gm^{\rot}}(\Gr_G; \C).
\end{equation*}
\end{lemma}
\begin{proof}
The affine Grassmannian $\Gr_G$ has connected components $\Gr^{\tau}_G$ labeled by elements $\tau \in \pi_1(G)$. The identity component $\Gr^1_G$ is $T^{\sc} \times \Gm^{\rot}$-equivariantly isomorphic to $\Gr_{G^{\sc}}$, which we have treated already in our main theorem.

In general, we can repeat the proof of the main theorem for all components at once, using the corresponding variant of cocharacter graphs from \S \ref{section: cocharacter graphs for reductive groups}. Alternatively, we can follow \cite{BF07} and bootstrap the current case as follows. For each $\tau \in \pi_1(G)$ choose a lift $\underline{\tau} \in X_{\bullet}$. Translation by $\tau$ then identifies the identity component $\Gr^1_G$ with $\Gr^{\tau}_G$. This isomorphism is twisted-equivariant for the $T^{\sc} \times \Gm^{\rot}$-action; the twisting is given by the conjugation with $t^\tau$. Under this twisting, we identify the copies of $T^{\sc} \times \Gm^{\rot} \times T \sslash W$ compatibly with the cocharacter graphs inside them as in \S \ref{subsubsection: cocharacter graphs for reductive groups}. The conclusion follows.
\end{proof}






\begin{remark}
The connected components of $\Gr_G$, labeled by the fundamental group $\pi_1(G)$, are all abstractly isomorphic; such isomorphisms are realized by translating with elements of the loop group $\lop G$, for example by an element corresponding to some $\lambda \in X_{\bullet}$. Each of these components is isomorphic to the affine Grassmannian $\Gr_{G^{\sc}}$ of $G^{\sc}$.

However, we would like to emphasize the following caveat: these components are \textit{not} isomorphic equivariantly, not even with respect to the extended torus $T \times \Gm^{\rot}$. Indeed, given two components $\Gr^{\sigma}_G$, $\Gr^{\tau}_G$ for some $\sigma$, $\tau \in \pi_1(G)$, their equivariant $K$-theory rings are usually different as $\O(T \times \Gm^{\rot})$-algebras.

Indeed, as we explained, the fiber of the $K$-theory ring over a point $(x, \zeta) \in T \times \Gm^{\rot}$ is given by the cohomology of fixed-points. Now, take $(1, \zeta)$ for generic $\zeta$, and consider the affine Grassmannian of $PGL_2$. The fixed points of $(1, \zeta)$ are given by a union of partial flag varieties. In the even component, the fixed-point scheme contains a unique connected component isomorphic to a point. In the odd component, the fixed-point scheme has only components given by partial flag varieties of positive dimension. The cohomology rings of the fixed-points are not the same.

The same situation occurs for cohomology in \cite[\S 3, proof of Theorem 1.(b)]{BF07}.
\end{remark}






\subsubsection{Equivariance with respect to actual torus.}
Above, we have always considered equivariance with respect to the maximal torus $T^{\sc}$ of the simply connected $G^{\sc}$. However, one can well ask about a description with respect to the actual maximal torus $T \leq G$. 

We first assume $G$ is semisimple, giving isogenies $T^{\sc} \to T \to T^{\ad}$. To simplify notation, we take only the identity component $\Gr_{G^{\sc}}$ of $\Gr_G$. 
\begin{lemma}\label{lemma: equivariant K-theory as deformation to normal cone -- isogeneous torus}
Assume $G$ is semisimple. We have 
\begin{equation*}
    \widehat{\O}(\Bl^{\circ}_{\widetilde{\eZ}}(T \times \Gm^{\rot} \times T^{\sc}\sslash W ))
    =\widehat{\O}(\Bl^{\circ}_{\widetilde{\eZ}}(\tfrac{T^{\sc} \times T^{\sc} \sslash W}{\pi_1(G)} \times \Gm^{\rot}))
    =K^{\top, 0}_{T \times \Gm^{\rot}}(\Gr_{G^{\sc}}; \C).
\end{equation*}
\end{lemma}
\begin{proof}
By Iwahori pavings, the equivariant $K$-theory of the affine Grassmannian is equivariantly formal. Since the action of $T^{\sc}$ factors through $T$, we know that $K^{\top, 0}_{T^{\sc}}(\Gr_G; \k)$ is the pullback of $K^{\top, 0}_{T \times \Gm^{\rot}}(\Gr_G; \k)$ along the $\pi_1(G)$-torsor $T^{\sc} \to T$. Conversely, the $T$-equivariant $K$-theory is given by descending the $T^{\sc}$-equivariant one, in other words by taking the quotient by the $\pi_1(G)$-action.

Now, on the inclusion $\Gamma \hookrightarrow T^{\sc} \times \Gm \times T^{\sc}$, this action is given by \S \ref{section: action of the fundamental group}. On this level, the quotient by $\pi_1(G)$ is \eqref{equation: descended Gamma eq 2}. Since the affine blowup construction and its Schubert completion is compatible with quotients by finite groups, we deduce the result.    
\end{proof}

\begin{remark}
Note the apparent break of symmetry between the two actions. This is due to our desire to explicitly see the equivariant base $T$ as the first copy. If we wish to keep the symmetry, the appropriate description via $\tfrac{T^{\sc} \times T^{\sc}}{\pi_1(G)}$ comes from \eqref{equation: descended Gamma eq 1}.
\end{remark}


\begin{remark}\label{remark: equivariant K-theory as deformation to normal cone -- isogeneous torus}
Now assume $G$ is reductive. Then we have
\begin{equation*}
\widehat{\O}(\Bl^{\circ}_{\widetilde{\eZ}}(T \times \Gm^{\rot} \times T^{\sc}\sslash W ))
=K^{\top, 0}_{T \times \Gm^{\rot}}(\Gr_{G^{\sc}}; \C).    
\end{equation*}
Indeed, this follows by first applying Lemma \ref{lemma: equivariant K-theory as deformation to normal cone -- isogeneous torus} with respect to the the maximal torus $T^{\ad}$ of the semisimple group $G^{\ad}$, and then base-changing along $T \to T^{\ad}$.
\end{remark}